% ==============================================================================
% RESULTADOS
% ==============================================================================
% Organizado por objetivo (general, OE1, OE2, OE3), con un cuadro SMART al
% inicio de cada uno. Cada cifra proviene de un artefacto del repositorio;
% la fuente se indica en la nota de cada cuadro o figura.
%
% Fuentes principales (repositorio logSguarDian salvo indicación):
%   - docs/cybersecurity-objectives-compliance.md (estado al 2026-10-02).
%   - training/results/v11_test_results.json (lectura de prueba 2026-09-26).
%   - training/models/parity_report.json (vigente, rf_v11/if_v10).
%   - docs/results.md (A19, A24, F4.4).
%   - benchmarks/results/extractor-2026-10-09T23-34-*.json (extractor, PR #97;
%     reemplaza la cifra F1.8 de docs/results.md, que no tenía artefacto).
%   - docs/vulnerable-projects-tests/ y docs/limitations.md (terceros).
%   - logSguarDian-vulnerable-project: docs/config3b-results.md,
%     attack-sim/ua_fp_study_results.json, exclusion_report.json.
%
% Imágenes: copiadas de docs/cibersecurity-images/ y training/results/ a
% figuras/resultados/. Las capturas que muestran una generación de modelo o
% un estado anterior se rotulan como tales en su leyenda; ninguna se
% presenta como cifra vigente si no lo es.
%
% Se omitieron obj3-b (reporte de validación de una generación anterior,
% redundante con obj3-a y con el cuadro vigente) y obj3-f (misma Ronda 4
% que obj3-g, solo Configuración 2).
\ifdefined\modoensamble
\else
	\documentclass[11pt, letterpaper]{report}
\usepackage[spanish, es-nodecimaldot, es-noquoting]{babel}
\selectlanguage{spanish}
\usepackage[utf8]{inputenc}
\usepackage[T1]{fontenc}

% ===============================================================================
% El estudiante debe llenar sus datos en esta sección para que la plantilla los 
% auto-importe y genere automáticamente las páginas de portada y de firmas 
% autorizadas.
% ===============================================================================
% Datos del estudiante:
% -------------------------------------------------------------------------------
% Nombre completo
\def \nombreestudiante {Sebastián Huertas Gómez}
% Carné
\def \uvgcarne {22295}
% Facultad
\def \uvgfacultad {Ingeniería}
% Carrera
\def \uvgcarrera {Ingeniería en Ciencia de la Computación y Tecnologías de la Información}

% Datos del trabajo:
% -------------------------------------------------------------------------------
% Título completo (tomado literalmente de protocolo_sebastian_v1 (3).pdf, sección 1)
\def \titulotesis {logSguarDian: Desarrollo de una librería para la detección de amenazas de seguridad en Endpoints de aplicaciones}
% Mes y año de entrega (confirmado por el estudiante: noviembre 2026)
\def \mesentrega {noviembre}
\def \anoentrega {2026}
% Asesor (según protocolo_sebastian_v1 (3).pdf, portada)
\def \nombreasesor {Ing. Erick Francisco Marroquín Rodríguez}
% El catedrático del curso (Ing. Bacilio Alexander Bolaños Lima) figura como
% primer examinador en la hoja de firmas, definido más abajo.

% Datos del tribunal examinador:
% -------------------------------------------------------------------------------
% Nombre del primer examinador (segunda línea de la hoja de firmas)
\def \nombreprimerex {Ing. Ma. Bacilio Alexander Bolaños Lima}

% Capítulos pre-definidos
% -------------------------------------------------------------------------------
% Comentar las líneas de las secciones que desean omitirse, por defecto se
% se incluyen todas.
% Se descomentan conforme se redacta cada capítulo (todavía sin archivo de
% contenido: prefacio, alcance, anexos, glosario).
\def \CAPprefacio {Prefacio}
\def \CAPantecedentes {Antecedentes}
\def \CAPalcance {Alcance}
\def \CAPmetodologia {Metodología}
\def \CAPresultados {Resultados}
\def \CAPdiscusion {Análisis de resultados}
\def \CAPanexos {Anexos}
\def \CAPglosario {Glosario}

% Formato y estilo de la plantilla
% -------------------------------------------------------------------------------
% Modo impresión: Puede des-comentar la siguiente línea para generar un documento pdf sin la portada, para cuando se desee imprimir el documento para encuadernación
% NOTA: la plantilla original traía esta línea activa por defecto, lo cual
% suprime la portada verde (main.tex hace \let\CAPportada\undefined cuando
% \printver está definido). Se comenta para que la versión digital, con
% portada, sea la que se genera por defecto.
%\def \printver {Versión del documento para impresión}

% Portada: Puede cambiarse la imagen en la portada al cambiar el nombre del 
% archivo siguiente. NOTA: debe tener la suficiente resolución para cubrir el área
% designada
\def \imagenportada {plantilla/portadacit.jpg}

% Referencias: Puede des-comentar la siguiente línea para utilizar el formato de referencias APA
%\def \usarAPA {Usar formato APA}

% Párrafo: Puede comentar la siguiente línea si desea emplear un formato de 
% párrafo distinto al establecido por defecto
\def \parpordefecto {Formato de párrafo por defecto}

% Capítulos y secciones: Puede des-comentar la siguiente línea para establecer el
% formato de los capítulos y secciones bajo el estándar original de UVG para
% trabajos de graduación. Este incluye: capítulos con numeración romana, secciones
% con letras mayúsculas, sub-secciones con números y sub-sub-secciones con letras
% minúsculas
%\def \capsecuvg {Formato UVG para capítulos y secciones}
% ================================================================================
% En este archivo se colocan opciones adicionales para modificar el formato de la
% plantilla, para emplearse en otros tipos de documentos que no sean trabajos de
% graduación. Si usted está trabajando su tesis, NO modifique este archivo
% ================================================================================
% Capítulos pre-definidos
% --------------------------------------------------------------------------------
% Comentar las líneas de las secciones que desean omitirse, por defecto se 
% se incluyen todas.
\def \CAPportada {Portada}
\def \CAPcaratula {Caratula}
\def \CAPfirmas {Hoja de firmas}
\def \CAPindice {Índice general}
\def \CAPfiguras {Listado de figuras}
\def \CAPcuadros {Listado de cuadros}
\def \CAPresumen {Resumen}
\def \CAPabstract {Resumen}
\def \CAPintroduccion {Introducción}
\def \CAPobjetivos {Objetivos}
\def \CAPjustificacion {Justificación}
\def \CAPmarcoteorico {Marco teórico}
\def \CAPconclusiones {Conclusiones}
\def \CAPrecomendaciones {Recomendaciones}
\def \CAPbibliografia {Bibliografía}

\usepackage{xcolor}
\usepackage{amsfonts}
\usepackage{amsmath}
\usepackage{amssymb}
\usepackage{amsthm}
\usepackage{mathtools}
\usepackage{graphicx}
\usepackage{xfrac}
\usepackage{float}
\usepackage[hypertexnames=false]{hyperref}
\usepackage[font=small]{caption}
\usepackage{subcaption}
% Cuadros: array (columnas L{} alineadas a la izquierda) y longtable (cuadros
% que ocupan más de una página). Se cargan aquí, no en
% 2-paquetes_y_comandos_usuario.tex, para que los capítulos de Metodología y
% Resultados no dependan de ese archivo.
\usepackage{array}
\usepackage{longtable}
\newcolumntype{L}[1]{>{\raggedright\arraybackslash}p{#1}}
\renewcommand{\arraystretch}{1.3}
\setlength{\LTpre}{0.3cm}
\setlength{\LTpost}{0.1cm}
% Margen de estiramiento para que los párrafos con palabras largas no
% generen líneas que sobresalgan del margen.
\setlength{\emergencystretch}{3em}
\usepackage{xpatch}
\usepackage{emptypage}
% htt: permite dividir con guion el texto en \texttt (nombres de archivo,
% clases y campos), que de otro modo sobresale del margen derecho.
\usepackage[htt]{hyphenat}
\usepackage{fancyhdr}
\usepackage[backend=biber, style=ieee]{biblatex}
\ifdefined\usarAPA
    \usepackage[backend=biber, style=apa]{biblatex}
\fi
\addbibresource{m-bibliografia.bib}

\usepackage[percent]{overpic}
\usepackage{chngcntr}

\ifdefined\CAPglosario
	\usepackage[numberedsection]{glossaries}
	\makeglossaries
    \newglossaryentry{rasp}
{
    name={RASP},
    description={Autoprotección de aplicaciones en tiempo de ejecución (\textit{Runtime Application Self-Protection}): mecanismo de seguridad que se ejecuta dentro de la propia aplicación y analiza sus peticiones para detectar y bloquear ataques.}
}

\newglossaryentry{waf}
{
    name={WAF},
    description={Cortafuegos de aplicaciones web (\textit{Web Application Firewall}): dispositivo o servicio perimetral que filtra el tráfico HTTP antes de que llegue a la aplicación.}
}

\newglossaryentry{modsecurity}
{
    name={ModSecurity},
    description={Motor de WAF de código abierto que se usa con reglas externas; en este trabajo se configuró con OWASP CRS.}
}

\newglossaryentry{crs}
{
    name={OWASP CRS},
    description={Conjunto de reglas genéricas de OWASP (\textit{Core Rule Set}) para ModSecurity, con niveles de paranoia que controlan qué tan estrictas son.}
}

\newglossaryentry{owasp}
{
    name={OWASP Top 10},
    description={Lista de las diez categorías de riesgo más críticas para aplicaciones web, publicada por OWASP; este trabajo usa la edición 2021.}
}

\newglossaryentry{mitre}
{
    name={MITRE ATT\&CK},
    sort={MITRE ATT&CK},
    description={Base de conocimiento de tácticas y técnicas de atacantes, usada aquí para asociar cada vector de ataque a una técnica.}
}

\newglossaryentry{cwe}
{
    name={CWE},
    description={Enumeración común de debilidades de software (\textit{Common Weakness Enumeration}): catálogo de tipos de defectos que originan vulnerabilidades.}
}

\newglossaryentry{sqli}
{
    name={SQL Injection},
    description={Ataque que inserta sentencias SQL en una entrada de la aplicación para alterar las consultas que ejecuta la base de datos.}
}

\newglossaryentry{xss}
{
    name={XSS},
    description={\textit{Cross-Site Scripting}: ataque que inyecta código ejecutable en una página que otros usuarios verán.}
}

\newglossaryentry{pathtraversal}
{
    name={Path Traversal},
    description={Ataque que manipula rutas de archivo (por ejemplo con \texttt{../}) para leer archivos fuera del directorio previsto; incluye la inclusión local de archivos (LFI).}
}

\newglossaryentry{cmdi}
{
    name={Command Injection},
    description={Ataque que logra ejecutar comandos del sistema operativo a través de una entrada de la aplicación.}
}

\newglossaryentry{middleware}
{
    name={\textit{middleware}},
    sort={middleware},
    description={Función que se ejecuta entre la recepción de una petición y la lógica de negocio de la aplicación; en Express se registra con \texttt{app.use}.}
}

\newglossaryentry{rf}
{
    name={Random Forest},
    description={Algoritmo supervisado que combina muchos árboles de decisión; en la librería es el único modelo con autoridad para bloquear una petición.}
}

\newglossaryentry{iforest}
{
    name={Isolation Forest},
    description={Algoritmo no supervisado de detección de anomalías que aísla las observaciones atípicas con particiones aleatorias; en la librería solo registra anomalías y nunca bloquea.}
}

\newglossaryentry{onnx}
{
    name={ONNX},
    description={Formato abierto para serializar modelos de aprendizaje automático, que permite ejecutarlos fuera de Python.}
}

\newglossaryentry{workerthreads}
{
    name={\texttt{worker\_threads}},
    sort={worker_threads},
    description={Módulo de Node.js que ejecuta código en hilos independientes del hilo principal; la inferencia de los modelos se ejecuta ahí.}
}

\newglossaryentry{eventloop}
{
    name={\textit{Event Loop}},
    sort={Event Loop},
    description={Bucle de eventos de Node.js que atiende las peticiones en un único hilo; bloquearlo degrada a toda la aplicación.}
}

\newglossaryentry{failopen}
{
    name={\textit{fail-open}},
    sort={fail-open},
    description={Política por la cual, ante un error o un tiempo de espera del módulo de detección, la petición se deja pasar en lugar de bloquearse.}
}

\newglossaryentry{monitor}
{
    name={Modo \texttt{monitor}},
    sort={Modo monitor},
    description={Modo de operación en el que la librería registra las peticiones que bloquearía, pero no las bloquea.}
}

\newglossaryentry{block}
{
    name={Modo \texttt{block}},
    sort={Modo block},
    description={Modo de operación en el que la librería responde HTTP 403 a las peticiones que clasifica como ataque.}
}

\newglossaryentry{pesp}
{
    name={p95},
    description={Percentil 95 del tiempo de respuesta: el valor por debajo del cual queda el 95~\% de las peticiones.}
}

\newglossaryentry{deltap95}
{
    name={$\Delta p95$},
    sort={Delta p95},
    description={Diferencia entre el p95 con el middleware activo y el p95 de la línea base sin él; en este trabajo se mide en milisegundos.}
}

\newglossaryentry{funo}
{
    name={F1},
    description={Media armónica de la precisión y el recall; es la métrica principal de clasificación del trabajo.}
}

\newglossaryentry{auc}
{
    name={AUC-ROC},
    description={Área bajo la curva ROC; mide qué tan bien separa un modelo dos clases sin depender de un umbral.}
}

\newglossaryentry{wilson}
{
    name={Intervalo de Wilson},
    description={Intervalo de confianza para una proporción, adecuado para muestras pequeñas o proporciones cercanas a 0 o 1.}
}

\newglossaryentry{casidup}
{
    name={Casi-duplicado},
    description={Registro muy parecido a otro pero no idéntico; si queda repartido entre entrenamiento y prueba, puede inflar las métricas.}
}

\newglossaryentry{minhash}
{
    name={MinHash},
    description={Técnica para estimar rápidamente la similitud entre conjuntos, usada para detectar casi-duplicados.}
}

\newglossaryentry{fuga}
{
    name={Fuga de datos},
    description={Presencia en el conjunto de evaluación de información disponible durante el entrenamiento, lo que sobreestima el desempeño.}
}

\newglossaryentry{corpuslimpio}
{
    name={Corpus limpio},
    description={Conjunto de payloads de evaluación filtrado para excluir los que coinciden con el entrenamiento; en este trabajo es el de la Ronda~5.}
}

\newglossaryentry{useragent}
{
    name={\texttt{User-Agent}},
    sort={User-Agent},
    description={Cabecera HTTP con la que el cliente se identifica (navegador, cURL, entre otros); su representación en el corpus sesgó al modelo.}
}

\newglossaryentry{falsopositivo}
{
    name={Falso positivo},
    description={Petición legítima que el sistema clasifica o bloquea como ataque.}
}

\newglossaryentry{parmodelos}
{
    name={Par de modelos},
    description={Combinación de un archivo \texttt{rf.onnx} y uno \texttt{if.onnx} identificada por su suma SHA-256 (pares A, B y C).}
}

\newglossaryentry{tarball}
{
    name={\textit{tarball}},
    sort={tarball},
    description={Archivo comprimido \texttt{.tgz} de un paquete de npm, publicado o construido localmente.}
}

\newglossaryentry{kanban}
{
    name={Kanban},
    description={Metodología de gestión basada en un tablero de columnas que representan el estado de las tareas.}
}

\newglossaryentry{artillery}
{
    name={Artillery},
    description={Herramienta de pruebas de carga usada para medir la latencia con y sin el middleware.}
}

\newglossaryentry{sqlmap}
{
    name={sqlmap},
    description={Herramienta de código abierto que automatiza la detección y la explotación de inyecciones SQL; generó los payloads de SQL Injection del corpus limpio.}
}

\newglossaryentry{zap}
{
    name={OWASP ZAP},
    description={Herramienta de OWASP para pruebas de seguridad de aplicaciones web; se usó para confirmar el payload de XSS.}
}

\newglossaryentry{mcp}
{
    name={MCP},
    description={Protocolo de Contexto de Modelo (\textit{Model Context Protocol}): estándar para que un asistente de inteligencia artificial use herramientas externas.}
}

\fi

\usepackage[top=1in, left=1.5in, right=1in, bottom=1in]{geometry}
\usepackage[Sonny]{fncychap}

\graphicspath{ {figuras/} }
\definecolor{uvg-green}{RGB}{17,71,52}
\newcommand{\defaultparformat}[1]{
	{\setlength{\parskip}{2ex}
     \input{#1}}
}
\ifdefined\capsecuvg
	\renewcommand\thechapter{\Roman{chapter}}
    \renewcommand\thesection{\Alph{section}}
	\renewcommand\thesubsection{\arabic{subsection}}
    \renewcommand\thesubsubsection{\alph{subsubection}}
\fi
\counterwithout{figure}{chapter}
\counterwithout{table}{chapter}
\counterwithout{equation}{chapter}

\providecommand{\blankpage}{
\newpage
\thispagestyle{empty}
\mbox{}
\newpage
}

% ==============================================================================
% Paquetes y comandos adicionales del usuario
% ==============================================================================
% TikZ: requerido por las figuras del capítulo de Marco Teórico (pipeline de
% middleware, comparación RASP vs. control perimetral).
\usepackage{tikz}
\usetikzlibrary{arrows.meta, positioning, shapes.geometric, fit, backgrounds}


\pagestyle{headings}
\begin{document}

\ifdefined\parpordefecto
	\setlength{\parskip}{2ex}
\fi

\fi

\chapter{Resultados}
\label{cap:resultados}

Este capítulo presenta la evidencia obtenida mediante los procedimientos del
Capítulo~\ref{cap:metodologia}, organizada por objetivo: primero el objetivo general y después
los tres objetivos específicos. Cada sección inicia con un cuadro que descompone el objetivo según
los criterios SMART \cite{doran1981smart} (específico, medible, alcanzable, relevante y temporal) y asocia cada criterio
con la evidencia correspondiente. A continuación se presentan los datos, mediciones y capturas que
sustentan esa evidencia. Cuando un indicador tiene un criterio de aceptación, se indica si el valor
medido lo cumple; la interpretación de los resultados se desarrolla en el capítulo de Análisis de resultados.

Cada resultado indica el par de modelos ONNX con que se obtuvo, según el
Cuadro~\ref{tab:pares-modelos}: las métricas sobre el conjunto de prueba, la evaluación en OWASP
Juice Shop y las Configuraciones 2 y 3b de la medición unificada del 25 de octubre de 2026
corresponden al par~C, que es el contenido del paquete 0.1.1; el paquete 0.1.0, la
medición principal de cobertura sobre el corpus limpio y el estudio de \texttt{User-Agent}, al
par~A; las aplicaciones DVNA, node-api-goat y
dvws-node, con atribución probable al par~A, pues el modelo exacto de cada medición no consta en el
repositorio; y el benchmark de rendimiento y la repetición de la Configuración~2 del 9 de
septiembre, al par~B. Algunas
capturas reflejan generaciones anteriores del modelo o del proyecto y se incluyen como valores
iniciales para la comparación; su leyenda lo indica de forma explícita.

% ==============================================================================
\section{Objetivo general}
\label{sec:res-og}

El objetivo general establece: diseñar e implementar una librería npm que integre un pipeline de
ingeniería de características y un modelo híbrido de aprendizaje automático capaz de detectar y
clasificar amenazas conocidas y anomalías estadísticas en endpoints Node.js en tiempo de
ejecución, sin comprometer el rendimiento del servidor anfitrión. El Cuadro~\ref{tab:res-smart-og}
lo descompone según los criterios SMART.

\begin{table}[H]
\centering
\small
\begin{tabular}{|L{2.2cm}|L{5.2cm}|L{5.6cm}|}
\hline
\textbf{Criterio} & \textbf{Definición en el objetivo} & \textbf{Evidencia obtenida} \\ \hline
Específico & Librería npm para Express con pipeline de características y modelo híbrido (Random
Forest + Isolation Forest) & Paquete \texttt{logsguardian} publicado (0.1.0 y 0.1.1); extractor de 76
características; RF con autoridad de bloqueo e IF de registro de anomalías
(Sección~\ref{subsec:res-og-libreria}) \\ \hline
Medible & Paridad ONNX $<$0.1~\%; extracción $\leq$5~ms; $\Delta p95_{\mathrm{abs}}\leq$15~ms;
$\Delta$ memoria $\leq$150~MB (143.1~MiB); CPU $\leq$15~\% & Valores en los Cuadros~\ref{tab:res-paridad},
\ref{tab:res-extractor}, \ref{tab:res-latencia} y \ref{tab:res-memoria} \\ \hline
Alcanzable & Stack de código abierto, sin costo de licencias & Construido con Node.js,
\texttt{onnxruntime-node}, SQLite, Docker y GitHub Actions, sin licencias de pago \\ \hline
Relevante & Protección a nivel de aplicación para organizaciones sin equipo de seguridad &
Instalación con un solo comando (\texttt{npm i logsguardian}); validación en cuatro aplicaciones de
terceros (Sección~\ref{subsec:res-og-terceros}) \\ \hline
Temporal & Mayo a octubre de 2026 & Publicación en npm el 30 de agosto de 2026 (0.1.0) y el 11 de octubre de 2026 (0.1.1); mediciones finales entre el 23 de septiembre y el 11 de octubre de 2026 \\ \hline
\end{tabular}
\caption[Criterios SMART del objetivo general]{Criterios SMART del objetivo general y evidencia asociada}
\label{tab:res-smart-og}
\end{table}
{\footnotesize Nota. Elaboración propia.}

\subsection{Librería publicada y verificación funcional}
\label{subsec:res-og-libreria}

La Figura~\ref{fig:res-npm} muestra el paquete publicado en el registro público de npm, en su
versión 0.1.0, con licencia MIT, publicada el 30 de agosto de 2026. El 11 de octubre de 2026 se
publicó la versión 0.1.1, hoy la vigente, que incorpora los modelos \texttt{rf\_v11} e
\texttt{if\_v10} y la corrección del defecto de recodificación de la consulta
(Cuadro~\ref{tab:pares-modelos}). La publicación se realiza de forma automática desde GitHub
Actions al fusionar la rama \texttt{develop} en \texttt{main}, siempre que la versión local difiera
de la publicada.

\begin{figure}[H]
\centering
\includegraphics[width=0.92\textwidth]{resultados/obj0-a.png}
\caption[Paquete logsguardian publicado en npm]{Página del paquete \texttt{logsguardian} en el registro público de npm (versión 0.1.0, captura previa a la publicación de la 0.1.1)}
\label{fig:res-npm}
\end{figure}
{\footnotesize Nota. Captura de npmjs.com/package/logsguardian. Fecha de publicación según el registro de npm.}

La Figura~\ref{fig:res-tests} muestra la estructura del monorepo (paquetes \texttt{core},
\texttt{extractor} y \texttt{mlops}) y la ejecución de la suite de pruebas del extractor de
características: 11 suites y 74 pruebas aprobadas, entre ellas las de determinismo,
decodificación recursiva y Command Injection compuesto.

\begin{figure}[H]
\centering
\includegraphics[width=0.92\textwidth]{resultados/obj0-b.png}
\caption[Paquetes del monorepo y pruebas del extractor]{Paquetes del monorepo y resultado de la suite de pruebas del extractor de características}
\label{fig:res-tests}
\end{figure}
{\footnotesize Nota. Captura de terminal, rama \texttt{feat/goals\_ciber}, repositorio \texttt{logSguarDian}.}

La Figura~\ref{fig:res-cli} muestra la CLI de la librería ejecutada dentro del contenedor de la
aplicación de referencia: \texttt{attacks list} cataloga los eventos clasificados por tipo (294
SQLi, 181 Path Traversal, 90 XSS y 23 Command Injection en la base de eventos de esa sesión), y
\texttt{attacks inspect sqli} muestra las métricas del modelo, las características de mayor
importancia y ejemplos de payloads detectados. En esa captura, de \texttt{rf\_v3}, la característica
de mayor importancia global del Random Forest es la longitud de la cabecera \texttt{User-Agent}
(\texttt{ua\_length}, 0.0796). En el análisis SHAP \cite{lundberg2017unified,valenzuela2026ia} de \texttt{rf\_v11}, \texttt{ua\_length} es también
la característica de mayor contribución media en las clases legítima (0.0866) y SQLi (0.0734).

\begin{figure}[H]
\centering
\includegraphics[width=\textwidth]{resultados/obj3-d.png}
\caption[CLI de logsguardian: attacks list e inspect]{Salida de los comandos \texttt{attacks list} y \texttt{attacks inspect sqli} de la CLI de \texttt{logsguardian}}
\label{fig:res-cli}
\end{figure}
{\footnotesize Nota. Captura de terminal, 22 de septiembre de 2026. Las métricas que muestra
\texttt{attacks inspect} en esta captura corresponden a \texttt{rf\_v3}; el archivo de métricas se
actualizó después a \texttt{rf\_v11}, y la versión 0.1.1 del paquete ya lo incluye
(Cuadro~\ref{tab:res-rf}).}

La Figura~\ref{fig:res-e2e} muestra la suite de pruebas extremo a extremo, que envía por HTTP real
100 payloads por clase a una aplicación Express con el middleware y los modelos ONNX publicados.
Las ocho pruebas se aprobaron: 100~\% de bloqueo en SQLi, 98~\% en XSS, 99~\% en Path Traversal,
96~\% en Command Injection y 2~\% de falsos positivos sobre peticiones legítimas. Como se indicó en
la Sección~\ref{subsec:met-fuga}, este conjunto comparte payloads con el entrenamiento; por ello
se presenta como verificación del funcionamiento extremo a extremo de la librería y no como
medición de detección sobre payloads nuevos, que se reporta en la Sección~\ref{subsec:res-oe3-cobertura}.

\begin{figure}[H]
\centering
\includegraphics[width=0.62\textwidth]{resultados/obj3-e.png}
\caption[Suite de pruebas extremo a extremo]{Resultado de la suite de detección extremo a extremo (\texttt{pnpm run test:e2e})}
\label{fig:res-e2e}
\end{figure}
{\footnotesize Nota. Captura de terminal, repositorio \texttt{logSguarDian}.}

\subsection{Paridad de la serialización ONNX}

El Cuadro~\ref{tab:res-paridad} presenta la paridad entre las predicciones de los modelos en
Python y sus versiones ONNX, calculada sobre 1{,}000 muestras para los modelos vigentes. Las
diferencias máximas son del orden de $10^{-7}$, por debajo del criterio de 0.1~\%. La
Figura~\ref{fig:res-paridad} muestra el reporte de paridad y la prueba automatizada que compara
\texttt{onnxruntime-node} con Python, con diferencia máxima de 0 en ambos modelos en esa captura, previa al reentrenamiento.

\begin{table}[H]
\centering
\small
\begin{tabular}{|L{4.6cm}|L{3.2cm}|L{3.2cm}|}
\hline
\textbf{Indicador} & \textbf{Random Forest} & \textbf{Isolation Forest} \\ \hline
Versión del modelo & \texttt{rf\_v11} & \texttt{if\_v10} \\ \hline
Características de entrada & 69 & 63 \\ \hline
Diferencia máxima Python--ONNX & $1.01\times10^{-7}$ & $2.42\times10^{-7}$ \\ \hline
Muestras comparadas & 1{,}000 & 1{,}000 \\ \hline
Criterio ($<$0.1~\%) & Cumple & Cumple \\ \hline
\end{tabular}
\caption[Paridad entre modelos Python y ONNX]{Paridad entre las predicciones de los modelos en Python y en ONNX}
\label{tab:res-paridad}
\end{table}
{\footnotesize Nota. Elaboración propia con datos de \texttt{training/models/parity\_report.json}.}

\begin{figure}[H]
\centering
\includegraphics[width=0.72\textwidth]{resultados/obj3-k.png}
\caption[Reporte y prueba de paridad ONNX]{Reporte de paridad ONNX y prueba automatizada de paridad entre Node.js y Python}
\label{fig:res-paridad}
\end{figure}
{\footnotesize Nota. Captura de terminal, repositorio \texttt{logSguarDian}. La captura es previa al
reentrenamiento del 26 de septiembre de 2026; los valores vigentes son los del
Cuadro~\ref{tab:res-paridad}.}

\subsection{Rendimiento: extracción de características y latencia}

El Cuadro~\ref{tab:res-extractor} presenta el tiempo de extracción del vector de características
por petición en dos series: la extracción directa en el hilo principal y la ida y vuelta completa
a un \texttt{worker\_thread} que ejecuta el extractor, que corresponde a la forma en que el
Alcance define el criterio. Cada serie se midió de forma serial sobre 2{,}000 iteraciones por tipo
de petición, en tres corridas independientes. En tráfico mixto, el p95 fue de 0.0126~ms en la
serie directa y de 0.0273~ms en la ida y vuelta al hilo; el mayor p95 de cualquier tipo de petición
en la ida y vuelta fue de 0.0305~ms (XSS). Todos los valores están por debajo del criterio de 5~ms.

\begin{table}[H]
\centering
\small
\begin{tabular}{|L{2.9cm}|L{1.5cm}|L{1.5cm}|L{1.5cm}|L{1.5cm}|L{1.5cm}|L{1.5cm}|}
\hline
 & \multicolumn{3}{L{5.34cm}|}{\textbf{Directa, hilo principal (ms)}} & \multicolumn{3}{L{5.34cm}|}{\textbf{Ida y vuelta a \texttt{worker\_thread} (ms)}} \\ \hline
\textbf{Tipo de petición} & \textbf{p50} & \textbf{p95} & \textbf{p99} & \textbf{p50} & \textbf{p95} & \textbf{p99} \\ \hline
Legítima & 0.0103 & 0.0169 & 0.0409 & 0.0213 & 0.0304 & 0.0467 \\ \hline
SQLi & 0.0092 & 0.0120 & 0.0340 & 0.0210 & 0.0275 & 0.0405 \\ \hline
XSS & 0.0102 & 0.0135 & 0.0340 & 0.0226 & 0.0305 & 0.0430 \\ \hline
Path Traversal & 0.0076 & 0.0097 & 0.0302 & 0.0162 & 0.0264 & 0.0377 \\ \hline
Command Injection & 0.0090 & 0.0102 & 0.0314 & 0.0212 & 0.0283 & 0.0405 \\ \hline
\textbf{Mixto} & \textbf{0.0097} & \textbf{0.0126} & \textbf{0.0350} & \textbf{0.0206} & \textbf{0.0273} & \textbf{0.0413} \\ \hline
\end{tabular}
\caption[Tiempo de extracción de características]{Tiempo de extracción del vector de características por petición, directa y con ida y vuelta a un \texttt{worker\_thread}}
\label{tab:res-extractor}
\end{table}
{\footnotesize Nota. Elaboración propia con datos de \texttt{benchmarks/results/extractor-2026-10-09T23-34-*.json}
(repositorio \texttt{logSguarDian}, commit \texttt{7b036bb}). Mediana de tres corridas; Apple M5,
Node.js v25.6.0; 200 iteraciones de calentamiento descartadas por serie.}

El Cuadro~\ref{tab:res-latencia} presenta la latencia del servidor anfitrión medida en Linux
nativo (GitHub Actions, \texttt{ubuntu-latest}, 26 de septiembre de 2026), con cinco repeticiones
por condición y la mediana como valor reportado. En los escenarios normal y de ataque, el
incremento absoluto del p95 fue de 12.08~ms y 12.85~ms respectivamente, ambos por debajo del
criterio de 15~ms. En el escenario de volumen, el criterio de 15~ms no se cumple: el p95 con middleware fue de 6{,}559.3~ms y el
7.00~\% de las peticiones (1{,}824 de 26{,}048) terminó por tiempo de espera excedido; en las
demás combinaciones de escenario y condición no hubo tiempos de espera excedidos. Los tiempos de
espera se concentraron en los últimos intervalos de cada repetición.

\begin{table}[H]
\centering
\small
\begin{tabular}{|L{2.3cm}|L{1.6cm}|L{1.6cm}|L{1.7cm}|L{1.7cm}|L{2.6cm}|}
\hline
\textbf{Escenario} & \textbf{p95 base (ms)} & \textbf{p95 activo (ms)} & \textbf{$\Delta p95_{\mathrm{abs}}$ (ms)} & \textbf{$\Delta p95_{\mathrm{rel}}$} & \textbf{Tiempos de espera excedidos} \\ \hline
Normal & 2.49 & 14.57 & 12.08 & +485.1~\% & 0 \\ \hline
Ataque & 4.62 & 17.47 & 12.85 & +278.1~\% & 0 \\ \hline
Volumen & 63.84 & 6{,}559.30 & 6{,}495.46 & +10{,}174.6~\% & 1{,}824 / 26{,}048 (7.00~\%) \\ \hline
\multicolumn{3}{|L{6.3cm}|}{Criterio ($\Delta p95_{\mathrm{abs}}\leq$15~ms)} & \multicolumn{3}{L{6.8cm}|}{Cumple en normal y en ataque; no cumple en volumen} \\ \hline
\end{tabular}
\caption[Latencia p95 con y sin middleware]{Latencia p95 del servidor anfitrión con y sin middleware, Linux nativo}
\label{tab:res-latencia}
\end{table}
{\footnotesize Nota. Elaboración propia con datos del flujo \texttt{latency-benchmark.yml}
(repositorio \texttt{logSguarDian-vulnerable-project}). Mediana de cinco repeticiones; el p95 de
cada repetición es el promedio ponderado de los p95 por intervalo posteriores al calentamiento.
Corrida 36276690236 de GitHub Actions (2 vCPU, 7.8~GiB), par de modelos B; evidencia cruda en
\texttt{docs/evidence/latency-linux-2026-09-26.tgz}.}

Como valores iniciales, la Figura~\ref{fig:res-lat-hist} (paneles~\ref{fig:res-lat-hist-rel} y~\ref{fig:res-lat-hist-abs})
presenta las mediciones previas en Docker Desktop sobre macOS, para cuatro variantes de la
arquitectura de escritura y de espera del Isolation Forest. La variante publicada (híbrida)
registró un incremento absoluto de 7 a 9~ms y relativo de +142~\% a +178~\%.

\begin{figure}[H]
\centering
\begin{subfigure}[t]{0.52\textwidth}
\centering
\includegraphics[width=\textwidth]{resultados/obj3-h1.png}
\caption{Incremento relativo}
\label{fig:res-lat-hist-rel}
\end{subfigure}
\hfill
\begin{subfigure}[t]{0.46\textwidth}
\centering
\includegraphics[width=\textwidth]{resultados/obj3-h2.png}
\caption{Incremento absoluto}
\label{fig:res-lat-hist-abs}
\end{subfigure}
\caption[Latencia inicial en Docker Desktop por variante de arquitectura]{Incremento de latencia p95 por variante de arquitectura, mediciones iniciales en Docker Desktop sobre macOS}
\label{fig:res-lat-hist}
\end{figure}
{\footnotesize Nota. Capturas de \texttt{docs/vulnerable-app-evaluation/config2-latency-evaluation.md}.
Mediciones previas a la metodología en Linux nativo del Cuadro~\ref{tab:res-latencia}.}

\subsection{Memoria y CPU}

El Cuadro~\ref{tab:res-memoria} presenta la memoria y el uso de CPU del contenedor de la
aplicación en la corrida de Linux nativo, con las muestras asignadas a cada repetición y el periodo
de calentamiento excluido. La memoria se expresa en MiB, la unidad que reporta \texttt{docker
stats}. Sin middleware, el contenedor usó en promedio entre 63 y 83~MiB; con el middleware activo,
entre 395 y 459~MiB. El incremento pareado de memoria fue de 332.02~MiB bajo carga normal,
331.86~MiB bajo carga de ataque y 375.92~MiB en el escenario de volumen. El criterio del protocolo
es un incremento de 150~MB (143.1~MiB) como máximo, y el valor medido lo supera más de dos veces en
los tres escenarios. El incremento de CPU fue de 88.93 puntos porcentuales bajo carga normal
(+614~\%) y de 85.12 bajo carga de ataque (+530~\%), por encima del criterio de 15~\%. El incremento
de memoria fue prácticamente igual en los escenarios normal y de ataque.

\begin{table}[H]
\centering
\small
\begin{tabular}{|L{1.8cm}|L{1.8cm}|L{1.8cm}|L{2.0cm}|L{2.0cm}|L{2.3cm}|}
\hline
\textbf{Escenario} & \textbf{Memoria base (MiB)} & \textbf{Memoria activo (MiB)} & \textbf{$\Delta$ memoria (MiB)} & \textbf{CPU base / activo (\%)} & \textbf{$\Delta$ CPU (pp; relativo)} \\ \hline
Normal & 65.18 & 397.51 & +332.02 & 14.48 / 103.32 & +88.93 (+614~\%) \\ \hline
Ataque & 63.22 & 394.78 & +331.86 & 16.07 / 101.25 & +85.12 (+530~\%) \\ \hline
Volumen & 83.29 & 459.15 & +375.92 & 82.34 / 115.29 & +32.99 (+40~\%) \\ \hline
\multicolumn{3}{|L{6.24cm}|}{Criterio: $\Delta$ memoria $\leq$150~MB (143.1~MiB); $\Delta$ CPU $\leq$15~\%} & \multicolumn{3}{L{7.14cm}|}{No cumple en ninguno de los escenarios} \\ \hline
\end{tabular}
\caption[Memoria y CPU con y sin middleware]{Memoria y CPU del contenedor de la aplicación con y sin middleware, Linux nativo}
\label{tab:res-memoria}
\end{table}
{\footnotesize Nota. Elaboración propia con datos de \texttt{docs/evidence/resources\_summary.md}
(repositorio \texttt{logSguarDian-vulnerable-project}), ventana A: periodo medido de cada
repetición, primeros 15~s excluidos; media de cinco repeticiones por condición e incremento como
media de las diferencias pareadas. Corrida 36276690236, par de modelos B, 2 vCPU. La memoria es
la del contenedor completo, no solo la del proceso Node.js. Los valores de CPU superiores a 100~\%
corresponden a más de un núcleo.}

\subsection{Validación en aplicaciones de terceros}
\label{subsec:res-og-terceros}

El Cuadro~\ref{tab:res-terceros} presenta los resultados de integrar el paquete publicado en npm, entonces la versión 0.1.0 (atribución
probable al par~A; el modelo exacto de cada medición no consta en el repositorio), en modo \texttt{block} y con la configuración por defecto, en tres aplicaciones vulnerables de
terceros. En cada una se confirmó primero la explotabilidad sin protección. Las tasas de bloqueo en
ataques se obtuvieron con el conjunto de payloads de la suite extremo a extremo
(Sección~\ref{sec:met-validez}), y las tasas sobre tráfico legítimo con peticiones de cada
aplicación. En modo \texttt{block}, la librería bloqueó entre el 92.5~\% y el 100~\% de los ataques y una proporción igual
de alta del tráfico legítimo con \texttt{User-Agent} de navegador, aunque las muestras de tráfico legítimo difieren en tamaño y
diversidad: en node-api-goat, las 50 peticiones se dirigieron a una sola ruta
(\texttt{/hexToRgb}), y la medición de DVNA, con 5 peticiones, es preliminar.

\begin{table}[H]
\centering
\small
\begin{tabular}{|L{2.0cm}|L{2.4cm}|L{1.3cm}|L{1.3cm}|L{1.4cm}|L{1.3cm}|L{1.8cm}|}
\hline
\textbf{Aplicación} & \textbf{Arquitectura} & \textbf{SQLi} & \textbf{XSS} & \textbf{Path Trav.} & \textbf{CMDi} & \textbf{Legítimo bloqueado} \\ \hline
DVNA \cite{appsecco2026dvna} & Express, EJS, MySQL & 38/40 & 39/40 & No aplica & 37/40 & 5/5 (100~\%)\textsuperscript{a} \\ \hline
node-api-goat \cite{layro012026nodeapigoat} & API pura sin base de datos & No aplica & 49/50 & 50/50 & 48/50 & 50/50 (100~\%)\textsuperscript{b} \\ \hline
dvws-node \cite{snoopysecurity2026dvwsnode} & Express, MongoDB/MySQL, SOAP/GraphQL & 40/40 & 37/40 & 40/40 & 40/40 & 37/40 (92.5~\%) \\ \hline
\end{tabular}
\caption[Bloqueo en aplicaciones de terceros]{Ataques y peticiones legítimas bloqueados en aplicaciones vulnerables de terceros}
\label{tab:res-terceros}
\end{table}
{\footnotesize Nota. Elaboración propia con datos de \texttt{docs/vulnerable-projects-tests/}. "No
aplica": la aplicación no expone un endpoint vulnerable a ese vector (verificado en el código).
En node-api-goat, sin base de datos, la categoría SQLi no se evaluó como vulnerabilidad.
\textsuperscript{a}Medición preliminar, $n=5$. \textsuperscript{b}50 valores distintos enviados a una
sola ruta (\texttt{/hexToRgb}), todos con el \texttt{User-Agent} de Chrome 128; una comprobación con
cinco peticiones y el mismo \texttt{User-Agent} bloqueó 4 de 5. Commits base evaluados, clonados el 30 de septiembre de 2026 y modificados con los cambios de integración y de compatibilidad registrados en \texttt{docs/vulnerable-projects-tests/}, sin commitear: DVNA \cite{appsecco2026dvna}, commit \texttt{9ba473a} (8 de noviembre de
2023); node-api-goat \cite{layro012026nodeapigoat}, commit \texttt{868ee6c} (12 de agosto de 2024);
dvws-node \cite{snoopysecurity2026dvwsnode}, commit \texttt{eb6e288} (7 de septiembre de 2026).}

Dos resultados complementan el cuadro. En dvws-node, la tasa de bloqueo de tráfico legítimo fue
de 75~\% (30/40) con el \texttt{User-Agent} por defecto de Node.js y de 92.5~\% (37/40) con un
\texttt{User-Agent} de navegador. En OWASP Juice Shop, evaluada por el componente de IA en modo
\texttt{monitor} con una copia local del paquete que contiene el par~C (\texttt{25b407e6...} y
\texttt{5be28ee8...}), y no con el paquete del registro, el modelo habría bloqueado 303 de 312
peticiones legítimas con \texttt{User-Agent} de navegador (97.1~\%) y 123 de 312 con cURL
(39.4~\%). El Cuadro~\ref{tab:res-terceros-ua} reúne, para cada medición de bloqueo de tráfico legítimo, el
modelo, el modo y el \texttt{User-Agent}, porque las cifras no son comparables entre sí si se omite
alguna de esas tres condiciones. En modo \texttt{monitor} la librería no bloquea: la cifra es la
proporción de peticiones a las que el modelo asignó veredicto de bloqueo.

\begin{table}[H]
\centering
\small
\begin{tabular}{|L{2.6cm}|L{2.2cm}|L{1.8cm}|L{3.0cm}|L{3.0cm}|}
\hline
\textbf{Aplicación} & \textbf{Modelo} & \textbf{Modo} & \textbf{User-Agent} & \textbf{Legítimo con veredicto de bloqueo} \\ \hline
Juice Shop & Par~C & \texttt{monitor} & Navegador / cURL & 303/312 (97.1~\%) / 123/312 (39.4~\%) \\ \hline
dvws-node & No consta & \texttt{block} & Navegador / Node.js por defecto & 37/40 (92.5~\%) / 30/40 (75~\%) \\ \hline
node-api-goat & No consta & \texttt{block} & Chrome 128 & 50/50 en una ruta; 4/5 en una comprobación \\ \hline
DVNA & No consta & \texttt{block} & No consta & 5/5, preliminar \\ \hline
Aplicación de referencia (Artillery) & Par~B & \texttt{block} & Artillery & 0 bloqueos (0 de 6{,}300 por repetición) \\ \hline
Aplicación de referencia (navegadores) & Par~A & \texttt{block} & 15 combinaciones de navegadores reales & 13/15 \\ \hline
\end{tabular}
\caption[Bloqueo de tráfico legítimo por condición]{Bloqueo de tráfico legítimo según aplicación, modelo, modo y \texttt{User-Agent}}
\label{tab:res-terceros-ua}
\end{table}
{\footnotesize Nota. Elaboración propia con datos de \texttt{docs/vulnerable-projects-tests/} y de las corridas de latencia y de User-Agent del proyecto. ``No consta'': el modelo exacto de la medición no se registró; la atribución probable de las tres aplicaciones de terceros es el par~A.}

En las tres aplicaciones de la tabla, la clase predicha de los ataques bloqueados
difirió con frecuencia de la clase real; por ejemplo, en DVNA, 27 de los 37 Command Injection
bloqueados se clasificaron como SQLi. Una quinta aplicación, OWASP NodeGoat, no se integró porque
su imagen base (Node.js 12) es incompatible con los requisitos de la librería (Node.js $\geq$20).

% ==============================================================================
\section{Objetivo específico 1: evaluación de los vectores de ataque}
\label{sec:res-oe1}

\begin{table}[H]
\centering
\small
\begin{tabular}{|L{2.2cm}|L{5.2cm}|L{5.6cm}|}
\hline
\textbf{Criterio} & \textbf{Definición en el objetivo} & \textbf{Evidencia obtenida} \\ \hline
Específico & Cuatro vectores intra-request: SQLi, XSS, Path Traversal/LFI y Command Injection &
Cinco clases canónicas (cuatro de ataque y una legítima) en la taxonomía del dataset \\ \hline
Medible & Mapeo de cada vector a OWASP Top 10 2021 y MITRE ATT\&CK & 4/4 vectores mapeados a
categoría OWASP, CWE y técnica ATT\&CK, verificados en fuente oficial \\ \hline
Alcanzable & Marcos públicos y gratuitos & Páginas oficiales de OWASP y MITRE \\ \hline
Relevante & Estructurar el dataset del OE2 y los criterios del OE3 & Las cinco clases definen la
etiqueta de M1 y las categorías de todas las métricas del OE3 \\ \hline
Temporal & Fase 1 (mayo de 2026) & Taxonomía versionada en junio de 2026 \\ \hline
\end{tabular}
\caption[Criterios SMART del objetivo específico 1]{Criterios SMART del objetivo específico 1 y evidencia asociada}
\label{tab:res-smart-oe1}
\end{table}
{\footnotesize Nota. Elaboración propia.}

El Cuadro~\ref{tab:res-smart-oe1} resume el objetivo. El mapeo de los cuatro vectores a CWE, OWASP
Top 10 2021 y MITRE ATT\&CK se presenta en el Cuadro~\ref{cuadro:owasp-mitre-cwe} del Marco
Teórico: SQL Injection, XSS y Command Injection corresponden a A03:2021 (Injection) y Path
Traversal/LFI a A01:2021 (Broken Access Control). El Cuadro~\ref{tab:res-oe1} presenta cómo cada
vector se materializó en el proyecto: su clase canónica, su origen en la fuente CAPEC del corpus, el
endpoint de la aplicación de referencia donde se evaluó y la herramienta que generó sus payloads de
evaluación.

\begin{table}[H]
\centering
\small
\begin{tabular}{|L{2.3cm}|L{2.0cm}|L{2.3cm}|L{3.4cm}|L{2.6cm}|}
\hline
\textbf{Vector} & \textbf{Clase canónica} & \textbf{Columna CAPEC} & \textbf{Endpoint de evaluación} & \textbf{Generación de payloads} \\ \hline
SQL Injection & \texttt{sqli} & 66 & \texttt{GET /posts/search?q=}, \texttt{POST /login} & sqlmap \\ \hline
XSS & \texttt{xss} & 242 (Code Injection) & \texttt{POST /posts}, \texttt{POST /profile} & Generador propio (semilla 20261025) \\ \hline
Path Traversal/LFI & \texttt{path\_traversal} & 126 & \texttt{GET /posts/:id/ attachment?file=} & Generador propio \\ \hline
Command Injection & \texttt{cmdi} & 88, 248 & \texttt{POST /admin/ping} & Generador propio \\ \hline
\end{tabular}
\caption[Materialización de los vectores en el proyecto]{Materialización de los cuatro vectores en la taxonomía del dataset y en el entorno de evaluación}
\label{tab:res-oe1}
\end{table}
{\footnotesize Nota. Elaboración propia con datos de \texttt{training/label\_map.yaml} y del
repositorio \texttt{logSguarDian-vulnerable-project}. La asignación de CAPEC-242 a \texttt{xss} es
una decisión de mapeo del dataset, no una equivalencia conceptual. El corpus de XSS lo genera
\texttt{generate\_xss.py} (400 candidatos, muestreo estratificado a 200, filtro de exclusión y
auditoría de fuga) con la semilla 20261025 el 25 de octubre de 2026; la semilla y la fecha difieren
de las de las otras tres categorías (20260923).}

La Figura~\ref{fig:res-labelmap} muestra la taxonomía de cinco clases y el mapeo de las columnas
CAPEC en el archivo versionado de etiquetas. La Figura~\ref{fig:res-featurespec} muestra un
fragmento de la especificación del vector de características, que define cada característica, su
fórmula y la categoría de ataque que discrimina; el vector completo tiene 76 características, de
las cuales el Random Forest consume 69 y el Isolation Forest 63.

\begin{figure}[H]
\centering
\includegraphics[width=0.68\textwidth]{resultados/obj1-a.png}
\caption[Taxonomía de clases en label\_map.yaml]{Taxonomía canónica de cinco clases y mapeo de columnas CAPEC en \texttt{label\_map.yaml}}
\label{fig:res-labelmap}
\end{figure}
{\footnotesize Nota. Captura de \texttt{training/label\_map.yaml}, repositorio \texttt{logSguarDian}.}

\begin{figure}[H]
\centering
\includegraphics[width=\textwidth]{resultados/obj1-b.png}
\caption[Fragmento de la especificación de características]{Fragmento de la especificación del vector de características}
\label{fig:res-featurespec}
\end{figure}
{\footnotesize Nota. Captura de \texttt{docs/feature-spec.md}, repositorio \texttt{logSguarDian}.}

% ==============================================================================
\section{Objetivo específico 2: dataset etiquetado}
\label{sec:res-oe2}

\begin{table}[H]
\centering
\small
\begin{tabular}{|L{2.2cm}|L{5.2cm}|L{5.6cm}|}
\hline
\textbf{Criterio} & \textbf{Definición en el objetivo} & \textbf{Evidencia obtenida} \\ \hline
Específico & Dataset de tráfico HTTP etiquetado con tráfico legítimo y las cuatro categorías &
Corpus M1 con cinco clases presentes en las tres particiones \\ \hline
Medible & Mínimo 100{,}000 muestras etiquetadas & 383{,}202 muestras (Cuadro~\ref{tab:res-particiones}) \\ \hline
Alcanzable & Fuentes públicas y conjuntos sintéticos & Siete fuentes públicas auditadas en el
corpus, más dos listas públicas de payloads y conjuntos propios
(Cuadro~\ref{tab:met-origenes}, Figura~\ref{fig:res-auditoria}) \\ \hline
Relevante & Base de entrenamiento y validación representativa & Partición estratificada por clase y
fuente con verificación de integridad \\ \hline
Temporal & Fase 1 (mayo de 2026) & Corpus inicial en junio de 2026; regenerado hasta septiembre de
2026 \\ \hline
\end{tabular}
\caption[Criterios SMART del objetivo específico 2]{Criterios SMART del objetivo específico 2 y evidencia asociada}
\label{tab:res-smart-oe2}
\end{table}
{\footnotesize Nota. Elaboración propia.}

El Cuadro~\ref{tab:res-smart-oe2} resume el objetivo. El Cuadro~\ref{tab:res-particiones} presenta
el tamaño de cada partición del corpus vigente: 383{,}202 muestras etiquetadas, 3.8 veces el mínimo
de 100{,}000 establecido. Las particiones de validación y de prueba tienen el mismo tamaño
(57{,}481 filas) por la proporción 70/15/15.

\begin{table}[H]
\centering
\small
\begin{tabular}{|L{4.0cm}|L{2.8cm}|L{2.8cm}|}
\hline
\textbf{Partición} & \textbf{Filas} & \textbf{Proporción} \\ \hline
Entrenamiento & 268{,}240 & 70~\% \\ \hline
Validación & 57{,}481 & 15~\% \\ \hline
Prueba & 57{,}481 & 15~\% \\ \hline
\textbf{Total} & \textbf{383{,}202} & \textbf{100~\%} \\ \hline
Criterio ($\geq$100{,}000) & \multicolumn{2}{L{6.0cm}|}{Cumple} \\ \hline
\end{tabular}
\caption[Tamaño de las particiones del corpus]{Tamaño de las particiones del corpus etiquetado}
\label{tab:res-particiones}
\end{table}
{\footnotesize Nota. Elaboración propia con datos de \texttt{training/results/v11\_test\_results.json}
y conteo directo de \texttt{training/splits/*.parquet}.}

El Cuadro~\ref{tab:res-clases} presenta la distribución por clase en el conjunto de prueba. El
corpus no está balanceado en conteos: SQLi representa el 59.3~\% de la partición y Command
Injection el 2.6~\%. El desbalance se trató con ponderación de clases en el entrenamiento. Como
referencia, antes de la deduplicación y el filtrado las fuentes crudas sumaban 974{,}718 filas
(57.8~\% legítimas, 30.3~\% SQLi, 7.1~\% Path Traversal, 3.8~\% XSS y 0.9~\% Command Injection).

\begin{table}[H]
\centering
\small
\begin{tabular}{|L{4.0cm}|L{2.8cm}|L{2.8cm}|}
\hline
\textbf{Clase} & \textbf{Filas} & \textbf{Proporción} \\ \hline
\texttt{sqli} & 34{,}107 & 59.3~\% \\ \hline
\texttt{benign} & 14{,}886 & 25.9~\% \\ \hline
\texttt{xss} & 4{,}491 & 7.8~\% \\ \hline
\texttt{path\_traversal} & 2{,}529 & 4.4~\% \\ \hline
\texttt{cmdi} & 1{,}468 & 2.6~\% \\ \hline
\textbf{Total} & \textbf{57{,}481} & \textbf{100~\%} \\ \hline
\end{tabular}
\caption[Distribución por clase del conjunto de prueba]{Distribución por clase del conjunto de prueba}
\label{tab:res-clases}
\end{table}
{\footnotesize Nota. Elaboración propia con datos de \texttt{training/results/v11\_test\_results.json}.}

La Figura~\ref{fig:res-train-dist} muestra la distribución por clase de la partición de
entrenamiento en la generación del corpus previa a su regeneración del 12 de septiembre de 2026,
con el mismo orden de magnitud por clase que la partición de prueba vigente.

\begin{figure}[H]
\centering
\includegraphics[width=0.62\textwidth]{resultados/obj2-b.png}
\caption[Distribución por clase de la partición de entrenamiento]{Carga de características y distribución por clase de la partición de entrenamiento}
\label{fig:res-train-dist}
\end{figure}
{\footnotesize Nota. Captura de notebook de entrenamiento, repositorio \texttt{logSguarDian}.
Generación previa a la regeneración del 12 de septiembre de 2026 (268{,}064 filas de entrenamiento y
57{,}443 de validación); el tamaño vigente es el del Cuadro~\ref{tab:res-particiones}.}

La Figura~\ref{fig:res-auditoria} muestra la tabla de auditoría de las nueve fuentes candidatas,
con su licencia verificada, DOI, formato, número de registros, categorías cubiertas y decisión de
inclusión. Se aprobaron ocho y se excluyó AIT-LDSv2; siete de las aprobadas aportan filas al corpus
(Cuadro~\ref{tab:met-origenes}). La Figura~\ref{fig:res-parsers}
muestra los nueve analizadores por fuente y el validador del esquema canónico, junto con los
documentos de metodología de etiquetado, deduplicación y muestreo.

\begin{figure}[H]
\centering
\includegraphics[width=\textwidth]{resultados/obj1-c.png}
\caption[Auditoría de fuentes del dataset]{Tabla de auditoría de las fuentes del dataset}
\label{fig:res-auditoria}
\end{figure}
{\footnotesize Nota. Captura de \texttt{docs/dataset-audit.md}, repositorio \texttt{logSguarDian}.}

\begin{figure}[H]
\centering
\includegraphics[width=\textwidth]{resultados/obj2-e.png}
\caption[Analizadores por fuente y documentación del etiquetado]{Analizadores por fuente y documentación de la metodología de etiquetado}
\label{fig:res-parsers}
\end{figure}
{\footnotesize Nota. Captura de terminal, repositorio \texttt{logSguarDian}.}

La Figura~\ref{fig:res-lock} muestra el historial de commits de bloqueo del conjunto de prueba, desde el inicial del 11 de junio
hasta el de \texttt{rf\_v11} del 30 de agosto de 2026. No hay un único bloqueo por generación: el valor del bloqueo cambió nueve
veces en total, la última el 23 de septiembre de 2026 (hash \texttt{7486c258...}), con el mismo
cambio repartido en tres commits, y la lectura del 26 de septiembre verificó ese hash antes de cualquier
inferencia.

\begin{figure}[H]
\centering
\includegraphics[width=0.78\textwidth]{resultados/obj2-f.png}
\caption[Historial de bloqueos del conjunto de prueba]{Historial de commits de bloqueo del conjunto de prueba}
\label{fig:res-lock}
\end{figure}
{\footnotesize Nota. Captura de terminal, repositorio \texttt{logSguarDian}. El hash mostrado
(\texttt{a3140ab3...}) corresponde al bloqueo del 30 de agosto de 2026, previo al bloqueo vigente.}

% ==============================================================================
\section{Objetivo específico 3: efectividad del sistema de detección}
\label{sec:res-oe3}

\begin{table}[H]
\centering
\small
\begin{tabular}{|L{2.2cm}|L{5.2cm}|L{5.6cm}|}
\hline
\textbf{Criterio} & \textbf{Definición en el objetivo} & \textbf{Evidencia obtenida} \\ \hline
Específico & Precisión, recall y F1 por categoría de ataque & Reporte de clasificación de
\texttt{rf\_v11} sobre el conjunto de prueba bloqueado (Cuadro~\ref{tab:res-rf}) \\ \hline
Medible & F1 $\geq$0.80 en al menos 3 de 4 categorías; cobertura $\geq$80~\% por categoría & F1
entre 0.9239 y 0.9956 (4/4); cobertura entre 92.5~\% y 100~\% en las 4 categorías con el paquete 0.1.1 (85.0~\% en Path Traversal con la 0.1.0) \\ \hline
Alcanzable & Conjunto de prueba bloqueado y entorno Docker reproducible & Lectura de prueba con
verificación de hash; corpus limpio reproducible con semilla fija \\ \hline
Relevante & Balance entre falsos negativos y falsos positivos sin analista dedicado & Falsos
positivos medidos en el conjunto de prueba, en la suite E2E y con navegadores reales \\ \hline
Temporal & Fases 4 y 5 (agosto a septiembre de 2026) & Lectura de prueba el 26 de septiembre;
corpus limpio el 24 de septiembre de 2026 \\ \hline
\end{tabular}
\caption[Criterios SMART del objetivo específico 3]{Criterios SMART del objetivo específico 3 y evidencia asociada}
\label{tab:res-smart-oe3}
\end{table}
{\footnotesize Nota. Elaboración propia.}

El Cuadro~\ref{tab:res-smart-oe3} resume el objetivo según los criterios SMART.

\subsection{Métricas de clasificación sobre el conjunto de prueba}

El Cuadro~\ref{tab:res-rf} presenta las métricas del Random Forest \texttt{rf\_v11} (par C) sobre
el conjunto de prueba bloqueado (57{,}481 filas), obtenidas en la lectura del 26 de septiembre de
2026, que verificó el hash del conjunto y las sumas de verificación de los binarios evaluados. Fue la tercera lectura
del conjunto de prueba para la generación \texttt{rf\_v11}; las dos anteriores se describen en la
Sección~\ref{subsec:met-f2}. Las métricas asignan a cada petición la clase de mayor probabilidad,
sin aplicar el umbral de bloqueo de 0.35. Las cuatro categorías de ataque superaron el
criterio de F1 $\geq$0.80. El F1 más bajo fue el de Command Injection (0.9239) y el más alto el de
SQLi (0.9956). La exactitud global fue de 0.9932 y el F1 macro de 0.9776. En la clase legítima, 23
de 14{,}886 peticiones (0.15~\%) se clasificaron como ataque.

\begin{table}[H]
\centering
\small
\begin{tabular}{|L{3.2cm}|L{1.8cm}|L{1.8cm}|L{1.8cm}|L{1.8cm}|L{1.6cm}|}
\hline
\textbf{Clase} & \textbf{Precisión} & \textbf{Recall} & \textbf{F1} & \textbf{AUC-ROC} & \textbf{Soporte} \\ \hline
\texttt{cmdi} & 0.9015 & 0.9475 & 0.9239 & 0.99901 & 1{,}468 \\ \hline
\texttt{path\_traversal} & 0.9853 & 0.9798 & 0.9826 & 0.99899 & 2{,}529 \\ \hline
\texttt{sqli} & 0.9955 & 0.9956 & 0.9956 & 0.99982 & 34{,}107 \\ \hline
\texttt{xss} & 0.9950 & 0.9804 & 0.9877 & 0.99930 & 4{,}491 \\ \hline
\texttt{benign} & 0.9984 & 0.9985 & 0.9984 & 0.99998 & 14{,}886 \\ \hline
\textbf{Macro} & 0.9751 & 0.9804 & \textbf{0.9776} & --- & 57{,}481 \\ \hline
\multicolumn{3}{|L{7.6cm}|}{Criterio (F1 $\geq$0.80 en $\geq$3 de 4)} & \multicolumn{3}{L{6.0cm}|}{Cumple (4 de 4)} \\ \hline
\end{tabular}
\caption[Métricas del Random Forest sobre el conjunto de prueba]{Métricas por clase del Random Forest (\texttt{rf\_v11}) sobre el conjunto de prueba bloqueado}
\label{tab:res-rf}
\end{table}
{\footnotesize Nota. Elaboración propia con datos de \texttt{training/results/v11\_test\_results.json}
(26 de septiembre de 2026). AUC-ROC en esquema uno contra el resto.}

Como valor inicial, la Figura~\ref{fig:res-rfv3} muestra las métricas de \texttt{rf\_v3} sobre el
conjunto de prueba de su generación (59{,}947 filas), con un F1 macro de 0.9682 y un F1 de 0.8954
en Command Injection. Entre \texttt{rf\_v3} y \texttt{rf\_v11}, el F1 de Command Injection pasó de
0.8954 a 0.9239, el de Path Traversal de 0.9671 a 0.9826, el de XSS de 0.9844 a 0.9877 y el de SQLi
de 0.9955 a 0.9956. Las particiones de ambas generaciones no son idénticas.

\begin{figure}[H]
\centering
\includegraphics[width=0.66\textwidth]{resultados/obj3-a.png}
\caption[Métricas iniciales de rf\_v3]{Métricas por clase de \texttt{rf\_v3} sobre el conjunto de prueba de su generación (valor inicial)}
\label{fig:res-rfv3}
\end{figure}
{\footnotesize Nota. Captura de \texttt{docs/decision-policy.md}, sección 2.1, repositorio \texttt{logSguarDian}.}

La Figura~\ref{fig:res-cmdi-evol} muestra una de las iteraciones de la Fase~5 sobre la categoría
de menor desempeño: la incorporación de 2{,}470 payloads reales y diversos de Command Injection en
\texttt{v9}, que elevó el F1 de esa clase en validación de 0.89--0.90 a 0.9369, después de que el
sobremuestreo sintético (SMOTE) produjera una mejora de solo 0.015.

\begin{figure}[H]
\centering
\includegraphics[width=0.7\textwidth]{resultados/obj2-d.png}
\caption[Evolución del F1 de Command Injection con datos reales diversos]{Evolución del desempeño en Command Injection al incorporar datos reales diversos (\texttt{v9})}
\label{fig:res-cmdi-evol}
\end{figure}
{\footnotesize Nota. Captura de \texttt{docs/limitations.md}, sección 1.1, repositorio \texttt{logSguarDian}.}

La Figura~\ref{fig:res-confusion} presenta la matriz de confusión del Random Forest sobre la
partición de validación. La confusión más frecuente fue la de SQLi clasificada como Command
Injection (106 casos), seguida de XSS clasificada como SQLi (56 casos).

\begin{figure}[H]
\centering
\includegraphics[width=0.74\textwidth]{resultados/rf_confusion_matrix.png}
\caption[Matriz de confusión del Random Forest]{Matriz de confusión del Random Forest sobre la partición de validación}
\label{fig:res-confusion}
\end{figure}
{\footnotesize Nota. Generada con scikit-learn \cite{pedregosa2011}; archivo
\texttt{training/results/rf\_confusion\_matrix.png}, 26 de septiembre de 2026.}

El Cuadro~\ref{tab:res-if} presenta las métricas del Isolation Forest sobre el mismo conjunto de
prueba, con el umbral fijado en sus metadatos (0.00421), seleccionado sobre validación. La tasa de detección de anomalías fue de
91.25~\% y la de falsos positivos de 5.46~\%; ambos criterios se cumplen. La
Figura~\ref{fig:res-ifcurve} muestra la curva de recall frente a la tasa de falsos positivos y el
punto de operación seleccionado.

\begin{table}[H]
\centering
\small
\begin{tabular}{|L{4.6cm}|L{2.6cm}|L{2.6cm}|L{2.2cm}|}
\hline
\textbf{Indicador} & \textbf{Valor} & \textbf{Criterio} & \textbf{Estado} \\ \hline
Detección de anomalías (recall) & 91.25~\% (38{,}868 / 42{,}595) & $\geq$50~\% & Cumple \\ \hline
Tasa de falsos positivos & 5.46~\% (813 / 14{,}886) & $\leq$10~\% & Cumple \\ \hline
Precisión & 97.95~\% & --- & --- \\ \hline
Recall por clase: \texttt{cmdi} / \texttt{path\_traversal} & 95.64~\% / 99.68~\% & --- & --- \\ \hline
Recall por clase: \texttt{sqli} / \texttt{xss} & 89.54~\% / 98.02~\% & --- & --- \\ \hline
\end{tabular}
\caption[Métricas del Isolation Forest sobre el conjunto de prueba]{Métricas del Isolation Forest (\texttt{if\_v10}) sobre el conjunto de prueba bloqueado}
\label{tab:res-if}
\end{table}
{\footnotesize Nota. Elaboración propia con datos de \texttt{training/results/v11\_test\_results.json}.}

\begin{figure}[H]
\centering
\includegraphics[width=0.8\textwidth]{resultados/if_recall_fp_curve.png}
\caption[Curva de recall frente a falsos positivos del Isolation Forest]{Recall frente a tasa de falsos positivos del Isolation Forest y punto de operación seleccionado}
\label{fig:res-ifcurve}
\end{figure}
{\footnotesize Nota. Archivo \texttt{training/results/if\_recall\_fp\_curve.png}, 26 de septiembre de 2026.}

\subsection{Cobertura sobre el corpus limpio}
\label{subsec:res-oe3-cobertura}

El Cuadro~\ref{tab:res-corpus} presenta la construcción del corpus limpio de la Ronda~5 y el
Cuadro~\ref{tab:res-cobertura} la cobertura obtenida con la Configuración~2. El corpus de la
Ronda~5 original contenía 601 peticiones: 200 por cada una de las categorías SQLi, Path Traversal y
Command Injection, y un payload de XSS confirmado con OWASP ZAP. El 25 de octubre de 2026 se agregó
un corpus propio de 200 payloads de XSS, con lo que el corpus de cobertura tiene 800 peticiones.
El filtro de exclusión descartó 3 candidatos de Path Traversal por
coincidencia exacta con una fuente de entrenamiento; en SQLi y Command Injection no hubo descartes.
Con el paquete 0.1.0 (par~A), la cobertura fue de 100~\% en SQLi, 91.5~\% en Command Injection
y 85.0~\% en Path Traversal, con intervalos de Wilson \cite{wilson1927probable} al 95~\% de 98.1--100.0~\%, 86.8--94.6~\% y
79.4--89.3~\%, respectivamente. Con el tarball de validación 22550ea (par~C) y el mismo corpus, la
cobertura fue de 99.5~\%, 96.0~\% y 95.0~\%, con intervalos de 97.2--99.9~\%, 92.3--98.0~\% y
91.0--97.3~\%. Las tres categorías superan el criterio de 80~\% con ambos pares, pero con el par~A
el límite inferior del intervalo de Path Traversal queda por debajo de ese umbral. Con el paquete
0.1.1, publicado después, instalado desde el registro de npm y medido con el ejecutor corregido, la
cobertura fue de 99.5~\% en SQLi, 100.0~\% en Path Traversal y 92.5~\% en Command Injection, con
intervalos de 97.2--99.9~\%, 98.1--100.0~\% y 88.0--95.4~\%. La medición del par~C no es una
evaluación estrictamente ciega, y la del paquete 0.1.1 se hizo sobre el corpus que reveló los
defectos de medición (Sección~\ref{sec:met-validez}). Cada medición corresponde a una sola corrida, salvo la de XSS, que se midió en tres corridas idénticas.
Dos peticiones se contaron como no detectadas por fallos y no por el veredicto del modelo:
\texttt{cmdi\_94} terminó por tiempo de espera en la corrida del par~A y \texttt{sqli\_75} devolvió
HTTP 500 en la del par~C. En XSS, el \textit{Active Scan} de ZAP, dirigido por su API REST y con
sesión autenticada, produjo 538 alertas, de las cuales una sola es de XSS: \texttt{POST /profile},
campo \texttt{bio}, con el payload \texttt{</div><scrIpt>alert(1);</scRipt><div>} y riesgo alto.
Los campos \texttt{content} y \texttt{title} de \texttt{POST /posts} se escanearon con la misma
política y no resultaron vulnerables. Aunque la regla de ZAP se llama \textit{Reflected}, la
vulnerabilidad es de XSS almacenado, porque \texttt{bio} se guarda y se vuelve a mostrar en cada
vista del perfil. Ese payload fue bloqueado por ambos pares de modelos (1 de 1), y el 25 de octubre de 2026 se
reconfirmó que el servidor lo almacena sin escapar (HTTP 200) sin protección y que la librería lo
bloquea (HTTP 403); la ejecución en un navegador no se verificó. Los
dos resultados responden preguntas distintas: ZAP (n=1) mide la explotabilidad confirmada, y el
corpus propio (n=200), la cobertura de detección.

La cobertura de XSS sobre el corpus propio, con el paquete 0.1.1 del registro de npm, fue de 200
de 200 payloads bloqueados (100.0~\%; intervalo de Wilson al 95~\%: 98.1--100.0~\%), sin variación
entre las tres corridas ni entre los tres transportes (64 de 64 en
\texttt{application/x-www-form-urlencoded}, 68 de 68 en JSON y 68 de 68 en
\texttt{multipart/form-data}); sin protección se bloquearon 0 de 200. De los 200 payloads, 64 se
dirigen al campo vulnerable \texttt{bio}, y los 64 se bloquearon (94.3--100.0~\%). Los otros 136 se
dirigen a \texttt{content} y \texttt{title}, campos sin vulnerabilidad confirmada: sin protección la
aplicación los rechaza con HTTP 400 porque la ruta exige ambos campos, y con la librería reciben
HTTP 403, de modo que miden la detección y no la explotabilidad. El filtro de exclusión de XSS se
aplicó después del muestreo y descartó 0 candidatos, y la auditoría posterior de fuga encontró 3
payloads con coincidencia normalizada por subcadena con el entrenamiento (Cuadro~\ref{tab:res-corpus});
sin ellos, la librería bloquea 197 de 197.

Los fallos de Path Traversal se concentraban en las categorías de codificación. El desglose por
estilo de codificación, calculado con \texttt{training/audit/breakdown\_config2.py} sobre los
resultados de la Configuración~2 con el ejecutor anterior, muestra que el paquete 0.1.0 (par~A) no
bloqueó 30 payloads: 21 de los 43 con codificación URL o mixta, 7 de los 23 con doble codificación
y 2 de los 19 con UTF-8 sobrelargo. Con el modelo vigente (par~C) los fallos eran 10: 9 de los 23
con doble codificación y 1 de los 19 con UTF-8 sobrelargo.

Esos fallos tenían dos causas independientes, y ninguna era una limitación de detección. La primera
era un defecto del middleware, que reconstruía la cadena de consulta a partir de \texttt{req.query},
ya decodificada una vez por Express, y le añadía una capa de codificación. La segunda era un defecto
del ejecutor de pruebas, que volvía a codificar los payloads de Path Traversal y Command Injection
que el generador ya había codificado. El Cuadro~\ref{tab:res-pt-causas} muestra que ninguna
corrección bastó por sí sola: con el paquete 0.1.1, que incluye la corrección del middleware, y el
ejecutor anterior, los 10 fallos de Path Traversal fueron los mismos que con el tarball sin esa
corrección (payload por payload), y desaparecieron solo al corregir también el ejecutor. Con ambas
correcciones, el paquete 0.1.1, instalado desde el registro de npm, bloqueó 200 de 200 payloads de
Path Traversal (intervalo de Wilson al 95~\%: 98.1--100.0~\%).

\begin{table}[H]
\centering
\small
\begin{tabular}{|L{5.4cm}|L{2.4cm}|L{2.8cm}|L{2.4cm}|}
\hline
\textbf{Paquete y ejecutor} & \textbf{Fallos de Path Traversal} & \textbf{Fallos de Command Injection} & \textbf{Fallos de SQLi} \\ \hline
0.1.0 (par~A), ejecutor anterior & 30 & 17 & 0 \\ \hline
Tarball 22550ea (par~C), ejecutor anterior & 10 & 8 & 1 \\ \hline
0.1.1 (par~C), ejecutor anterior & 10 & 8 & 1 \\ \hline
0.1.1 (par~C), ejecutor corregido & 0 & 15 (1 por tiempo de espera) & 1 \\ \hline
\end{tabular}
\caption[Fallos por paquete y versión del ejecutor]{Payloads no bloqueados de cada categoría, por paquete y por versión del ejecutor de pruebas}
\label{tab:res-pt-causas}
\end{table}
{\footnotesize Nota. Elaboración propia con datos de los archivos de resultados de la Configuración~2
(repositorio \texttt{logSguarDian-vulnerable-project}). No se midió el paquete 0.1.0 con el
ejecutor corregido, de modo que la mejora de 85.0~\% a 100.0~\% en Path Traversal no puede
repartirse entre el cambio de modelo, la corrección del middleware y la del ejecutor.}

Con el ejecutor corregido, la cobertura de Command Injection pasó de 96.0~\% a 92.5~\% (185 de 200;
intervalo de Wilson al 95~\%: 88.0--95.4~\%). Esto no es un empeoramiento de la detección. Los
generadores producen 105 de los 200 payloads de Command Injection con separadores ya codificados
(\texttt{\%26}, \texttt{\%7c}, \texttt{\%0a}, \texttt{\%3B}), y el ejecutor anterior los codificaba
una segunda vez. De esos, 10 se bloqueaban en su forma corrompida y no se bloquean al enviarse
correctamente (todos reciben HTTP 200); otros 3 hicieron el recorrido contrario, y 1 terminó por
tiempo de espera. El resultado neto es de 7 payloads bloqueados menos. El 96.0~\% estaba inflado
por ese defecto, y el 92.5~\% es la primera medición correcta de esta técnica de evasión: la
librería no detecta 10 payloads de Command Injection con separadores codificados, un límite real de
la detección.

La corrección del ejecutor afecta también a las demás mediciones de Path Traversal y Command
Injection de este capítulo con el ejecutor anterior (paquete 0.1.0 y tarball de validación), que no
se repitieron; las cuatro configuraciones se midieron de nuevo con el ejecutor corregido el 25 de
octubre de 2026 (Cuadro~\ref{tab:res-ronda5}).

El filtro de exclusión descarta únicamente las coincidencias exactas del campo completo (y, en
Command Injection, las de plantilla), por lo que no garantiza la ausencia de solapamiento por
forma normalizada. La medición de solapamiento del trabajo de graduación de Diego Valenzuela \cite{valenzuela2026ia}
(sección VIII.3.6), con \texttt{training/leakage\_check.py} (commit \texttt{22550ea}) sobre las
particiones de entrenamiento, validación y prueba, da los siguientes resultados. Con igualdad
normalizada del campo completo (decodificación de URL, minúsculas y dígitos sustituidos), 4 de los
200 payloads de SQLi coinciden con alguna partición (4 con entrenamiento, 1 con validación y 2
con prueba) y ninguno de Path Traversal ni de Command Injection. Con la regla de subcadena, un
criterio permisivo que su autor califica de límite superior y no de medida de solapamiento, Path
Traversal alcanza 7 coincidencias y Command Injection 2. La auditoría de subcadena repetida el 25 de
octubre de 2026 sobre los cuatro corpus da 4 coincidencias en SQLi, 7 en Path Traversal, 2 en
Command Injection y 3 en XSS, de 200 payloads cada uno. La cobertura de SQLi reportada incluye,
por lo tanto, esos 4 payloads. Las cifras anteriores a esta ronda no son comparables: el corpus
de la Ronda~4 tiene 200 de 200 payloads de Command Injection y 7 de 113 de XSS que coinciden por
igualdad normalizada, y ninguno de SQLi ni de Path Traversal; por subcadena, 26 de SQLi y 13 de
XSS.

\begin{table}[H]
\centering
\small
\begin{tabular}{|L{2.7cm}|L{4.2cm}|L{1.8cm}|L{1.8cm}|L{1.4cm}|}
\hline
\textbf{Categoría} & \textbf{Fuente} & \textbf{Candidatos} & \textbf{Descartados} & \textbf{Finales} \\ \hline
SQLi & sqlmap v1.10.9 & 1{,}995 & 0 & 200 \\ \hline
Path Traversal & Generador propio & 400 & 3 & 200 \\ \hline
Command Injection & Generador propio (no derivado de SecLists) & 400 & 0 & 200 \\ \hline
XSS (hallazgo de ZAP) & OWASP ZAP (\textit{Active Scan}, API REST) & 1 alerta confirmada & 0 & 1 \\ \hline
XSS (cobertura) & Generador propio (\texttt{generate\_xss.py}) & 400 & 0\textsuperscript{a} & 200 \\ \hline
\end{tabular}
\caption[Construcción del corpus limpio]{Construcción del corpus limpio de evaluación (Ronda~5, semilla 20260923; XSS, semilla 20261025)}
\label{tab:res-corpus}
\end{table}
{\footnotesize Nota. Elaboración propia con datos de \texttt{docs/config3b-results.md},
\texttt{attack-sim/exclusion\_report.json}, \texttt{attack-sim/exclusion\_report\_xss\_20261025.json}
y \texttt{attack-sim/leakage\_report\_xss\_20261025.json} (repositorio
\texttt{logSguarDian-vulnerable-project}). \textsuperscript{a}Los 200 payloads finales de XSS son
el resultado del muestreo estratificado de los 400 candidatos, no del filtro. El filtro de exclusión
se aplicó después del muestreo, en orden inverso al de las otras categorías, y descartó 0; la
auditoría posterior de fuga encontró 3 payloads (\texttt{xss\_118}, \texttt{xss\_254} y
\texttt{xss\_342}) con coincidencia normalizada por subcadena con las particiones de entrenamiento y
de prueba, que no se corrigieron.}

\begin{table}[H]
\centering
\small
\begin{tabular}{|L{2.2cm}|L{1.9cm}|L{2.0cm}|L{2.0cm}|L{2.2cm}|L{1.9cm}|}
\hline
\textbf{Categoría} & \textbf{Ronda 4 (par A)} & \textbf{Ronda 5, 0.1.0 (par A)\textsuperscript{a}} & \textbf{Ronda 5, tarball 22550ea (par C)\textsuperscript{a}} & \textbf{Ronda 5, 0.1.1 (par C), ejecutor corregido\textsuperscript{b}} & \textbf{Criterio ($\geq$80~\%)} \\ \hline
SQLi & 76/77 (98.7~\%) & 200/200 (100.0~\%) & 199/200 (99.5~\%) & 199/200 (99.5~\%) & Cumple \\ \hline
Path Traversal & 197/200 (98.5~\%) & 170/200 (85.0~\%) & 190/200 (95.0~\%) & 200/200 (100.0~\%) & Cumple \\ \hline
Command Injection & 200/200 (100.0~\%) & 183/200 (91.5~\%) & 192/200 (96.0~\%) & 185/200 (92.5~\%) & Cumple \\ \hline
XSS (corpus propio, n=200) & 110/113 (97.3~\%) & No medido\textsuperscript{c} & No medido\textsuperscript{c} & 200/200 (100.0~\%)\textsuperscript{c} & Cumple \\ \hline
XSS (ZAP, n=1) & --- & 1/1 (n=1, corpus anterior, no comparable) & 1/1 (n=1, corpus anterior, no comparable) & 1/1 (n=1) & No aplica; explotabilidad confirmada en un caso \\ \hline
\end{tabular}
\caption[Cobertura por categoría con logsguardian]{Cobertura por categoría con \textit{logSguarDian} activo (Configuración~2)}
\label{tab:res-cobertura}
\end{table}
{\footnotesize Nota. Elaboración propia con datos de \texttt{docs/config3b-results.md} y
\texttt{attack-sim/results\_config2-develop-tarball-20261009.json} (repositorio
\texttt{logSguarDian-vulnerable-project}). El corpus de la Ronda~4 tiene fuga confirmada en Command Injection (200/200 payloads coinciden por plantilla con
una fuente de entrenamiento). \textsuperscript{a}Medido con la versión del ejecutor anterior a la
corrección descrita en el texto, que añadía una capa de codificación a los payloads de Path
Traversal y Command Injection ya codificados. \textsuperscript{b}Paquete instalado desde el registro de
npm (\texttt{attack-sim/results\_config2-published-0.1.1-harness-fixed-20261010.json}); es la única
columna medida con el ejecutor corregido. Las columnas de la Ronda~5 con el ejecutor anterior no
son comparables con ella. \textsuperscript{c}La cobertura de XSS sobre el corpus propio se midió el
25 de octubre de 2026 con el paquete 0.1.1 del registro (archivo SHA-256 \texttt{abf8d5c4...}) y el
emisor \texttt{send\_xss\_corpus.py}, en la corrida 1 de tres idénticas
(\texttt{attack-sim/results\_combined\_config2\_run1\_20261025.json}); intervalo de Wilson al 95~\%:
98.1--100.0~\%. No se midió con el paquete 0.1.0 ni con el tarball de validación. El caso de ZAP se
midió con un corpus anterior de un solo payload y no es comparable con el de 200.}

\subsection{Comparación con un WAF y defensa en profundidad}

Las Rondas 1 a 3, con un corpus pequeño construido a mano, produjeron tres resultados: ModSecurity
con OWASP CRS en nivel de paranoia 1 bloqueó los 9 ataques originales y las 14 variantes de evasión
(0 de 14 llegaron a la aplicación); con el WAF en modo de solo detección, \textit{logSguarDian} bloqueó inicialmente 16 de los 23
payloads; los 7 no detectados eran de XSS enviados como \texttt{multipart/form-data}, que la
librería no inspeccionaba porque el analizador de ese tipo de cuerpo se montaba por ruta, después
del middleware. Tras montar un analizador global antes del middleware, \textit{logSguarDian}
bloqueó los 23 (23/23), y esa misma corrección evidenció una vulnerabilidad de XSS almacenado en la
aplicación de referencia.

La Figura~\ref{fig:res-ronda4} presenta la Ronda~4, que comparó las cuatro configuraciones sobre
590 payloads de SecLists. La configuración sin protección no bloqueó ningún ataque (0/590), el WAF
solo bloqueó 550/590 (93.2~\%), \textit{logSguarDian} solo 583/590 (98.8~\%) y ambas capas 587/590
(99.5~\%). De los 40 ataques que el WAF no bloqueó, \textit{logSguarDian} bloqueó 37 (92.5~\%;
intervalo de Wilson al 95~\%: 80.1--97.4~\%). Esta ronda se midió el 15 de agosto de 2026 con el
tarball local de la 0.1.0 (par~A) y con el ejecutor anterior a la corrección, y la Configuración~3a se midió
en una rama distinta de la aplicación, de modo que no es un resultado de la 0.1.1. Ninguno de esos 40 ataques es de Command Injection (fueron 11 de SQLi, 3 de
XSS y 26 de Path Traversal); por igualdad normalizada no se detectó coincidencia con el
entrenamiento en ninguno de ellos, aunque por la regla permisiva de subcadena coinciden 7 de los 11
de SQLi y 1 de los 3 de XSS. La corrección del ejecutor no altera las peticiones de esas
categorías, pero la ronda no se repitió con el ejecutor corregido. Los resultados
de Command Injection de esta ronda están afectados por la fuga descrita en la nota del
Cuadro~\ref{tab:res-cobertura}; la comparación se repitió después con el corpus limpio de la Ronda~5 (Cuadro~\ref{tab:res-ronda5}). La
Configuración~2 se repitió sobre los mismos 590 payloads el 9 de septiembre de 2026 con el par~B y
también bloqueó 583 de 590, aunque 10 payloads cambiaron de estado entre ambas corridas. En esa
repetición, los 200 payloads de Command Injection se registraron con la clase \texttt{sqli}; la
confusión afecta la clase reportada en el registro, no la decisión de bloqueo.

\begin{figure}[H]
\centering
\includegraphics[width=0.88\textwidth]{resultados/obj3-g.png}
\caption[Comparación de configuraciones en la Ronda 4]{Ataques bloqueados por configuración y análisis de defensa en profundidad (Ronda~4)}
\label{fig:res-ronda4}
\end{figure}
{\footnotesize Nota. Captura de \texttt{docs/config3b-results.md}, repositorio
\texttt{logSguarDian-vulnerable-project}.}

El Cuadro~\ref{tab:res-ronda5} repite la comparación con el corpus limpio de la Ronda~5, que no
tiene la fuga de la Ronda~4, ampliado con el corpus propio de XSS (800 peticiones). A diferencia de
la medición del 10 de octubre de 2026, que se descartó, las cuatro configuraciones se ejecutaron de
nuevo el 25 de octubre de 2026 sobre el mismo árbol de la aplicación, con el mismo
\texttt{Dockerfile} y el ejecutor corregido, y con el paquete 0.1.1 en las Configuraciones 2 y 3b
(Procedimiento~4.1). La Configuración~1 bloqueó 0 de 800 payloads, lo que confirma que el corpus y
el ejecutor miden lo que declaran. En total, la librería sola (2) bloqueó 784 de 800 payloads
(98.0~\%), el WAF solo (3a) 749 (93.6~\%) y ambas capas (3b) 797 (99.6~\%).

Por categoría, el WAF solo bloqueó el 99.5~\% de SQLi, igual que la librería, el 98.5~\% de Path
Traversal, el 82.0~\% de Command Injection (intervalo de Wilson al 95~\%: 76.1--86.7~\%) y el
94.5~\% de XSS. En Path Traversal no hay ventaja del WAF: con el paquete 0.1.1 la librería bloqueó
200 de 200 y el WAF 197 de 200, una diferencia de tres payloads con intervalos de Wilson que se
solapan (98.1--100.0~\% y 95.7--99.5~\%). En Command Injection la librería bloqueó el 92.5~\%, el
WAF el 82.0~\% y la combinación el 98.5~\%. El 68.0~\% que se reportó para el WAF en la medición del
10 de octubre se reprodujo al volver a ejecutar el ejecutor anterior sobre este mismo árbol (136 de
200, con el mismo veredicto en los 600 payloads de esa medición), de modo que la diferencia con el
82.0~\% se debe al ejecutor, que volvía a codificar los separadores ya codificados, y no a la
aplicación. Aun con el ejecutor corregido, el WAF bloquea menos los 105 payloads de Command
Injection con separadores ya codificados (80 de 105, 76.2~\%) que los 95 sin codificar (84 de 95,
88.4~\%); la diferencia es de 12.2 puntos porcentuales (prueba exacta de Fisher bilateral,
$p = 0.028$; intervalo de confianza de Wald al 95~\% de la diferencia, de 1.8 a 22.6 puntos;
Cuadro~\ref{tab:anx-cmdi-codificacion}). Los dos subconjuntos no están controlados por otras
características del payload, por lo que la diferencia no se atribuye a la codificación como causa. En XSS, los 11 payloads que el WAF no bloqueó
usan doble codificación URL y se envían como JSON o \texttt{multipart/form-data}; la librería sola
los bloqueó todos.

De los 51 ataques que el WAF solo no bloqueó (1 de SQLi, 3 de Path Traversal, 36 de Command
Injection y 11 de XSS), la Configuración~3b bloqueó 48 y la Configuración~2, 44 (48 con las cabeceras del proxy). Para el WAF, los
ataques no bloqueados pasan de 68 a 40 en las tres categorías comparables con la medición anterior
(de 64 a 36 en Command Injection), y la reejecución del ejecutor anterior sobre el mismo árbol
muestra que esa reducción se debe al ejecutor. No se compara con las mediciones anteriores la
Configuración~3b ni la 2, porque provienen de otras versiones de la librería: la Configuración~2
anterior usó el par~A y de la 3b anterior no consta el paquete. Frente a la Configuración~2 sin las cabeceras del proxy, la combinación bloqueó 13 payloads
más (1 de SQLi y 12 de Command Injection), pero esa diferencia mezcla dos condiciones: la librería
dentro de la 3b recibe las cabeceras que agrega el proxy y la Configuración~2 no. La descomposición
es 13 = 10 + 3: las cabeceras elevan en 10 los payloads que bloquea la librería sola (794 de 800
en la Configuración~2 con cabeceras, frente a 784), y en condiciones equiparadas el WAF suma 3 de
800 payloads a la librería (\texttt{cmdi\_9}, \texttt{cmdi\_51} y \texttt{cmdi\_94}, de Command
Injection y bloqueados por la 3a). La librería suma 48 payloads sobre el WAF solo
(Sección~\ref{anx:sensibilidad}). La complementariedad es asimétrica: en Path Traversal, SQLi y XSS
la librería sola, con cabeceras, bloquea todo lo que bloquea la combinación y el WAF no agrega
nada, y solo en Command Injection cada capa bloquea payloads que la otra no (3 el WAF y 33 la
librería). Tres
payloads de Command Injection (\texttt{cmdi\_6}, \texttt{cmdi\_184} y \texttt{cmdi\_185}) no los
bloquea ninguna configuración. Cuatro payloads de la misma categoría (\texttt{cmdi\_112},
\texttt{cmdi\_134}, \texttt{cmdi\_31} y \texttt{cmdi\_38}) los bloquea la 3b y no los bloquean ni la
3a ni la Configuración~2: con un servidor de captura se comprobó que el cuerpo y el objetivo de la
petición llegan idénticos por ambas rutas y que el proxy agrega cuatro cabeceras \texttt{X-*}, y al
enviar los cuatro payloads a la Configuración~2 con esas cabeceras el paquete 0.1.1 los bloqueó en
las tres repeticiones, y sin ellas no bloqueó ninguno. La diferencia queda explicada por el proxy
(el recuento de cabeceras fuera de la lista estándar que lee el extractor pasa de 0 a 4), de modo
que la Configuración~2 y la librería dentro de la 3b no son la misma medición
(Sección~\ref{sec:met-validez}).

\begin{table}[H]
\centering
\small
\setlength{\tabcolsep}{4pt}
\begin{tabular}{|L{2.2cm}|L{1.7cm}|L{2.1cm}||L{2.2cm}||L{2.2cm}|L{2.2cm}|}
\hline
\textbf{Categoría} & \textbf{Config.~1 (sin protección)} & \textbf{Config.~2 (solo \textit{logSguarDian}, 0.1.1)} & \textbf{Config.~2 con cabeceras (post hoc)} & \textbf{Config.~3a (solo WAF)} & \textbf{Config.~3b (ambas capas, 0.1.1)} \\ \hline
SQLi & 0/200 (0.0~\%) & 199/200 (99.5~\%) & 200/200 (100.0~\%) & 199/200 (99.5~\%) & 200/200 (100.0~\%) \\ \hline
Path Traversal & 0/200 (0.0~\%) & 200/200 (100.0~\%) & 200/200 (100.0~\%) & 197/200 (98.5~\%) & 200/200 (100.0~\%) \\ \hline
Command Injection & 0/200 (0.0~\%) & 185/200 (92.5~\%) & 194/200 (97.0~\%) & 164/200 (82.0~\%) & 197/200 (98.5~\%) \\ \hline
XSS & 0/200 (0.0~\%) & 200/200 (100.0~\%) & 200/200 (100.0~\%) & 189/200 (94.5~\%) & 200/200 (100.0~\%) \\ \hline
\textbf{Total} & 0/800 (0.0~\%) & 784/800 (98.0~\%) & 794/800 (99.2~\%) & 749/800 (93.6~\%) & 797/800 (99.6~\%) \\ \hline
\end{tabular}
\caption[Cuatro configuraciones sobre el corpus limpio de la Ronda 5 con XSS]{Payloads bloqueados por configuración sobre el corpus limpio de la Ronda~5 ampliado con XSS (800 peticiones)}
\label{tab:res-ronda5}
\end{table}
{\footnotesize Nota. Elaboración propia con datos de
\texttt{attack-sim/results\_combined\_config\{1,2,3a,3b\}\_run1\_20261025.json} (repositorio
\texttt{logSguarDian-vulnerable-project}); conteos recalculados a partir de los archivos de
resultados. Medición del 25 de octubre de 2026, corrida 1 de tres idénticas, con el mismo árbol de
la aplicación y el mismo \texttt{Dockerfile} en las cuatro configuraciones, el ejecutor corregido y
el corpus combinado de 800 peticiones (SHA-256 \texttt{f1fb0b0a...}). Configuraciones 2 y 3b con el
paquete 0.1.1 del registro de npm (archivo SHA-256 \texttt{abf8d5c4...}; modelos \texttt{25b407e6...}
y \texttt{5be28ee8...}, par~C; ver Sección~\ref{subsec:met-f2}). Intervalos de Wilson al 95~\% en
SQLi, Path Traversal, Command Injection y XSS: Configuración~2, 97.2--99.9, 98.1--100.0,
88.0--95.4 y 98.1--100.0~\%; Configuración~2 con cabeceras, 98.1--100.0, 98.1--100.0, 93.6--98.6 y
98.1--100.0~\%; Configuración~3a, 97.2--99.9, 95.7--99.5, 76.1--86.7 y 90.4--96.9~\%;
Configuración~3b, 98.1--100.0, 98.1--100.0, 95.7--99.5 y 98.1--100.0~\%. Las mediciones anteriores
de la Configuración~2 con el paquete 0.1.0 (Cuadro~\ref{tab:res-cobertura}) no se incluyen ni se
comparan con la Configuración~3a. La columna Config.~2 se midió sin las cuatro cabeceras
\texttt{X-*} que agrega el proxy, por lo que describe un despliegue directo, sin proxy; se fijó
antes de conocer el efecto de esas cabeceras y es la cifra principal de la librería. La columna
Config.~2 con cabeceras es una sensibilidad añadida después de observar que la 3b bloqueaba
payloads que la Configuración~2 no (análisis post hoc); se midió en tres corridas idénticas y es la
comparable con la 3a y la 3b, porque la librería dentro de la 3b recibe esas cabeceras. En
condiciones equiparadas (3b frente a esa columna) el aporte del WAF sobre la librería es de 3
payloads (Sección~\ref{anx:sensibilidad}).}

\subsection{Falsos positivos con navegadores reales}

El Cuadro~\ref{tab:res-ua} presenta el estudio de falsos positivos con cadenas
\texttt{User-Agent} de cinco navegadores reales sobre tres rutas sin consulta ni cuerpo de la
aplicación de referencia. De las 15 combinaciones legítimas, 13 (86.7~\%) se bloquearon. Solo
Firefox en macOS evitó el bloqueo en dos rutas (\texttt{/login} y \texttt{/posts}). Como contraste,
la suite E2E registró 2 de 100 peticiones legítimas bloqueadas (Figura~\ref{fig:res-e2e}).

\begin{table}[H]
\centering
\small
\begin{tabular}{|L{3.6cm}|L{2.2cm}|L{2.2cm}|L{2.2cm}|L{2.0cm}|}
\hline
\textbf{\texttt{User-Agent}} & \textbf{\texttt{/login}} & \textbf{\texttt{/posts}} & \textbf{\texttt{/profile}} & \textbf{Bloqueadas} \\ \hline
Chrome (macOS) & 403 & 403 & 403 & 3/3 \\ \hline
Chrome (Windows) & 403 & 403 & 403 & 3/3 \\ \hline
Safari (macOS) & 403 & 403 & 403 & 3/3 \\ \hline
Safari (iOS) & 403 & 403 & 403 & 3/3 \\ \hline
Firefox (macOS) & 200 & 200 & 403 & 1/3 \\ \hline
\textbf{Total} & 4/5 & 4/5 & 5/5 & \textbf{13/15} \\ \hline
\end{tabular}
\caption[Falsos positivos con User-Agent de navegadores reales]{Respuesta del middleware a peticiones legítimas con \texttt{User-Agent} de navegadores reales}
\label{tab:res-ua}
\end{table}
{\footnotesize Nota. Elaboración propia con datos de \texttt{attack-sim/ua\_fp\_study\_results.json}
(repositorio \texttt{logSguarDian-vulnerable-project}). \texttt{/login} sin sesión; \texttt{/posts} y
\texttt{/profile} con sesión iniciada. Ejecutado con el par~A; atribución por la fecha del commit, ya
que el archivo de resultados no registra el modelo.}

% ==============================================================================
\section{Resumen de indicadores}
\label{sec:res-resumen}

El Cuadro~\ref{tab:res-resumen} consolida los indicadores definidos en el
Cuadro~\ref{tab:met-metricas} de la Metodología, con el valor medido y su estado frente al
criterio.

\begin{longtable}{|L{3.5cm}|L{1.2cm}|L{2.7cm}|L{3.0cm}|L{2.1cm}|}
\hline
\textbf{Indicador} & \textbf{Obj.} & \textbf{Criterio} & \textbf{Valor medido} & \textbf{Estado} \\ \hline
\endfirsthead
\multicolumn{5}{l}{\footnotesize\textit{(continuación)}} \\ \hline
\textbf{Indicador} & \textbf{Obj.} & \textbf{Criterio} & \textbf{Valor medido} & \textbf{Estado} \\ \hline
\endhead
\multicolumn{5}{r}{\footnotesize\textit{Continúa en la página siguiente}} \\
\endfoot
\endlastfoot
Vectores mapeados a OWASP, CWE y ATT\&CK & OE1 & 4/4 & 4/4 & Cumple \\ \hline
Tamaño del corpus & OE2 & $\geq$100{,}000 & 383{,}202 & Cumple \\ \hline
Clases representadas por partición & OE2 & 5/5 & 5/5 & Cumple \\ \hline
F1 por categoría & OE3 & $\geq$0.80 en $\geq$3 de 4 & 0.9239--0.9956 (4/4) & Cumple \\ \hline
AUC-ROC por categoría & OE3 & Referencia & 0.99899--0.99982 & --- \\ \hline
Detección del IF & OE3 & $\geq$50~\% & 91.25~\% & Cumple \\ \hline
Falsos positivos del IF & OE3 & $\leq$10~\% & 5.46~\% & Cumple \\ \hline
Cobertura SQLi / PT / CMDi (paquete 0.1.0, ejecutor anterior) & OE3 & $\geq$80~\% & 100 / 85.0 / 91.5~\% (PT, IC 95~\%: 79.4--89.3~\%) & Cumple; no robusto en PT \\ \hline
Cobertura SQLi / PT / CMDi (modelo vigente, tarball) & OE3 & $\geq$80~\% & 99.5 / 95.0 / 96.0~\% (PT, IC 95~\%: 91.0--97.3~\%) & Cumple \\ \hline
Cobertura SQLi / PT / CMDi (paquete 0.1.1, ejecutor corregido) & OE3 & $\geq$80~\% & 99.5 / 100.0 / 92.5~\% (PT, IC 95~\%: 98.1--100.0~\%; CMDi, IC 95~\%: 88.0--95.4~\%) & Cumple \\ \hline
Cobertura XSS (corpus propio, paquete 0.1.1) & OE3 & $\geq$80~\% & 200/200 (100.0~\%; IC 95~\%: 98.1--100.0~\%); 64/64 en \texttt{bio} (IC 95~\%: 94.3--100.0~\%) & Cumple en cobertura (n=200) \\ \hline
Explotabilidad de XSS confirmada (ZAP) & OE3 & Descriptivo & 1 caso (\texttt{POST /profile}, \texttt{bio}; n=1) & --- \\ \hline
Falsos positivos con navegador & OE3 & Descriptivo & 13/15 combinaciones de \texttt{User-Agent} & --- \\ \hline
Paridad ONNX & OG & $<$0.1~\% & $\leq2.42\times10^{-7}$ & Cumple \\ \hline
Extracción de características (p95, ida y vuelta al hilo) & OG & $\leq$5~ms & 0.0273~ms & Cumple \\ \hline
$\Delta p95_{\mathrm{abs}}$ normal / ataque & OG & $\leq$15~ms & 12.08 / 12.85~ms & Cumple \\ \hline
$\Delta p95_{\mathrm{rel}}$ normal / ataque & OG & Complementario & +485.1 / +278.1~\% & --- \\ \hline
$\Delta p95_{\mathrm{abs}}$ volumen (con 7.00~\% de tiempos de espera) & OG & $\leq$15~ms (referencia) & 6{,}495.46~ms & No cumple \\ \hline
Memoria adicional bajo carga (normal / ataque) & OG & $\leq$150~MB (143.1~MiB) & +332.02 / +331.86~MiB & No cumple \\ \hline
CPU adicional bajo carga (normal / ataque) & OG & $\leq$15~\% & +614 / +530~\% & No cumple \\ \hline
Contribución en defensa en profundidad (Ronda~5, 25 de octubre de 2026) & OG & Descriptivo & La librería suma 48 de 800 payloads al WAF solo (la 3b bloqueó 48 de los 51 ataques que el WAF no bloquea); en condiciones equiparadas (3b frente a la Configuración~2 con cabeceras) el WAF suma 3 de 800 a la librería & --- \\ \hline
Bloqueo de legítimos en terceros (Cuadro~\ref{tab:res-terceros-ua}) & OG & Descriptivo & 92.5--100~\% (modo \texttt{block}); 97.1~\% veredicto en Juice Shop (modo \texttt{monitor}) & --- \\ \hline
\caption[Resumen de indicadores]{Resumen de indicadores: criterio, valor medido y estado}
\label{tab:res-resumen}
\end{longtable}
{\footnotesize Nota. Elaboración propia con datos de las secciones anteriores de este capítulo.}

\ifdefined\modoensamble
\else
	\section*{Referencias}
	\printbibliography[heading=none]
	\end{document}
\fi
